% Folio 01, analysis: the Übungstest playbook. Skript ch. 1-4 (Erwartungshorizont items 1-45).
% The Vorbereitung's 16 Kernfähigkeiten as 16 moves (skeleton, fresh worked example, traps, tasks),
% the theorem catalogue, the counterexamples, exam craft, a 14-day plan, and a key that gives
% answers and first moves only.
\documentclass{kaelearn}
\subject{Analysis I}
\folio{01}
\title{Sixteen\\Moves}
\shorttitle{Test playbook}
\subtitle{Everything the Übungstest can ask,\\and the move that answers it.}
\sources{Skript ch.~1--4 · Erwartungshorizont 1--45 · Vorbereitung Übungstest · UE~1--5}
\readtime{2 min view · 3 h folio · 14 days of drills}
\filed{2026-11-21}
\plate{plate.png}{Zeno's nave}
\kael{keeper}

\makeatletter
% tasks T01, T02, ... carry labels, so the key always matches
\newcounter{task}
\renewcommand\thetask{T\twodigit{\arabic{task}}}
\newcommand\task[1]{\par\smallskip\noindent\refstepcounter{task}\label{#1}%
  {\tek\fontsize{9.5}{10}\selectfont\color{klhi}\thetask}\hspace{7pt}}
\newenvironment{yourturn}{\begin{panel}{Your turn}{answers and first moves in the key}{}}{\end{panel}}
\newcommand\keyline[1]{\par\smallskip\noindent\hangindent=2.7em\hangafter=1%
  \makebox[2.7em][l]{\tek\fontsize{9.5}{10}\selectfont\color{klhi}\ref{#1}}}
% move furniture
\newcommand\when[1]{\par\smallskip\noindent\hangindent=3.2em\hangafter=1%
  \makebox[3.2em][l]{\color{kldim}\pix{When}}#1\par\smallskip}
\newcommand\trap[1]{\par\smallskip\noindent\hangindent=3.2em\hangafter=1%
  \makebox[3.2em][l]{\color{kldim}\pix{Trap}}#1\par}
% the two documents of an epsilon proof, side by side
\newenvironment{twodocs}{\par\smallskip\noindent\tabularx{\linewidth}{@{}X@{\hspace{7mm}}X@{}}%
  {\color{kldim}\pix{Scratch // backwards}} & {\color{kldim}\pix{Proof // forwards}}\\\kl@rule{2}}%
  {\endtabularx\par\smallskip}
% tables as rows of minipages, so a page may break between any two rows
\newlength\klcolw
\newcommand\kl@cell[2]{\begin{minipage}[t]{#1}\raggedright\strut#2\strut\end{minipage}}
\newcommand\kl@sep{\par\nopagebreak\vspace{1pt}\noindent\ledline{\linewidth}\par\vspace{2pt}}
% definitions and their negations
\newenvironment{negtable}{\par}{\par}
\newcommand\negrow[3]{\par\needspace{4\baselineskip}\noindent{\color{klhi}#1}\par\nopagebreak
  \noindent\makebox[11mm][l]{\color{kldim}\pix{says}}$#2$\par\nopagebreak
  \noindent\makebox[11mm][l]{\color{kldim}\pix{fails}}$#3$\kl@sep}
% two-column tables: #1 width of the left column, #2 and #3 the headers
\newenvironment{twocol}[3]{\par\needspace{6\baselineskip}\setlength\klcolw{#1}\row{\color{kldim}\pix{#2}}{\color{kldim}\pix{#3}}\nopagebreak}{\par}
\newcommand\row[2]{\noindent\kl@cell{\klcolw}{#1}\hspace{8pt}\kl@cell{\dimexpr\linewidth-\klcolw-8pt\relax}{#2}\kl@sep}
% the plan
\newenvironment{plan}{\par\pday{\color{kldim}\pix{Day}}{\color{kldim}\pix{Do}}{\color{kldim}\pix{Done when}}}{\par}
\newcommand\pday[3]{\noindent\kl@cell{8mm}{{\tek\color{klhi}#1}}\hspace{6pt}\kl@cell{78mm}{#2}\hspace{6pt}%
  \kl@cell{\dimexpr\linewidth-86mm-12pt\relax}{{\color{kldim}#3}}\kl@sep}
\renewenvironment{triggers}{\begin{panel}{Triggers}{if you see this, reach for that}{}\par
  \setlength\klcolw{58mm}\row{\color{kldim}\pix{You see}}{\color{kldim}\pix{Reach for}}}{\end{panel}}
\renewcommand\trigger[2]{\row{#1}{#2}}
\pretocmd{\section}{\needspace{14\baselineskip}}{}{}
\pretocmd{\subsection}{\needspace{8\baselineskip}}{}{}
\makeatother
% dictionary entries and the theorem catalogue
\newcommand\entry[2]{\par\needspace{3\baselineskip}\smallskip\noindent{\color{klhi}#1}\enspace\de{(#2)}\enspace}
\newcommand\eg{\enspace{\color{kldim}\pix{e.g.}}\enspace}
\newcommand\nonex{\enspace{\color{kldim}\pix{not}}\enspace}
\newcommand\satz[2]{\par\needspace{5\baselineskip}\medskip\noindent{\color{klhi}\textbf{#1}}\enspace\de{(#2)}\par\nopagebreak\noindent\ignorespaces}
\newcommand\idea{\par\nopagebreak\smallskip\noindent\hangindent=2.8em\hangafter=1\makebox[2.8em][l]{\color{kldim}\pix{Idea}}}
\newcommand\N{\mathbb N}
\newcommand\Q{\mathbb Q}
\newcommand\R{\mathbb R}
\newcommand\eps{\varepsilon}
\newcommand\cmd[1]{{\mono\fontsize{9}{11}\selectfont #1}}

\begin{document}
\makecover

% =====================================================================================
\section{Knowing is not solving}

\opening{Y}{ou can recite} the $\eps$-criterion. Ask you to prove $\frac{n^2+2n}{2n^2-5}\to\frac12$ with it, and the pen
stops. That is not a gap in knowledge. It is a gap in \emph{moves}: the small set of standard manoeuvres that turn a definition
into a proof. Nobody invents them during an exam. They are learned, the way a chess player learns openings instead of deriving
every game from the rules.
\aside{I have read every analysis exam ever set. They are the same forty ideas wearing different numbers.}

Your professor has been more generous than he may have meant to be. The Erwartungshorizont lists every item he can ask, and
it says what an answer to each kind looks like:
\begin{itemize}
  \item a \textbf{term} \de{(Begriff)}: the definition, an example, a counterexample;
  \item a \textbf{theorem} \de{(Satz)}: the statement with every hypothesis, and the proof or its idea;
  \item a \textbf{formula}: the formula, and the methods of computing with it from the exercises.
\end{itemize}
The Vorbereitung sheet names sixteen core skills \de{(Kernfähigkeiten)} up to the end of chapter~4. They are the sixteen moves
of this folio, in the same order, so the two documents lie side by side. Each move gets the question shapes that call for it,
a skeleton you can write half asleep, a worked example that is \emph{not} the one on the sheet, its traps, and problems for
you. Part~VII holds every theorem in the shape the test wants, Part~VIII every counterexample, Part~IX the craft of the exam
itself, Part~X your fourteen days.

One thing I will not do: hand you full solutions to the problems marked T. The key gives the answer and the first move. The
road from there to the proof is the exact skill the test measures, and it only grows when you walk it yourself.

\begin{twominute}
\begin{itemize}
  \item \textbf{The test.} Open questions on Skript chapters 1--4: Erwartungshorizont items 1--45, in the shapes of the
        Vorbereitung's sixteen core skills.
  \item \textbf{Three answer shapes.} Term: definition, example, counterexample. Theorem: every hypothesis, then the proof
        idea. Computation: the method from the exercises.
  \item \textbf{Every $\eps$-proof is written twice.} Scratch work runs backwards, from the goal $|a_n-a|<\eps$ to a condition
        on $n$. The proof runs forwards: choose $N$, then one chain of inequalities that ends in $<\eps$.
  \item \textbf{Estimates do the work.} Make the expression simpler and bigger until it reads $\frac Cn$ or $C\,|x-x_0|$.
        Crude is fine. Nobody grades the size of your $N$.
  \item \textbf{To refute, negate.} Flip every quantifier, then build one counterexample. For functions it is usually two
        sequences.
  \item \textbf{One axiom carries the course.} Archimedes, $\sqrt a$, nested intervals, monotone convergence,
        Bolzano--Weierstraß, Cauchy, the intermediate and extreme value theorems: the supremum axiom in costumes.
  \item \textbf{Skript tools only.} No derivatives, no L'Hospital, no $(1+\frac1n)^n\to e$. If you need a fact the Skript
        lacks, prove it in two lines or take another road.
\end{itemize}
\end{twominute}

\begin{guessfirst}
\begin{enumerate}
  \item A sequence has exactly one accumulation point \de{(Häufungspunkt)}. Must it converge?
  \item $x\mapsto x^2$ is continuous on $\R$. Is it uniformly continuous there? And $x\mapsto\sqrt x$ on $[0,\infty)$?
  \item Let $a_1=1$ and $a_{n+1}=2a_n$. If $a_n\to L$, then $L=2L$, so $L=0$. Where is the mistake?
\end{enumerate}
\end{guessfirst}

\begin{terms}
\term{supremum}{Supremum}{least upper bound; need not lie in the set}
\term{Archimedean property}{Archimedische Eigenschaft}{every real number is exceeded by some $n\in\N$}
\term{nested intervals}{Intervallschachtelung}{closed, bounded, each inside the last}
\term{accumulation point}{Häufungspunkt}{the limit of some subsequence}
\term{Cauchy sequence}{Cauchy-Folge}{late terms close to each other}
\term{absolutely convergent}{absolut konvergent}{$\sum|a_n|$ converges}
\term{uniformly continuous}{gleichmäßig stetig}{one $\delta$ serves every point}
\term{intermediate value theorem}{Zwischenwertsatz}{continuous on $[a,b]$: every value in between is hit}
\end{terms}

% =====================================================================================
\section{The engine}

\subsection{Quantifiers are a game}
\opening{E}{very definition} in this course is a game against an opponent. Read $\forall$ as \emph{he chooses} and $\exists$
as \emph{you choose}, strictly from left to right:
\[
  \forall\eps>0\ \ \exists N\in\N\ \ \forall n\ge N:\quad |a_n-a|<\eps .
\]
He names an $\eps$, as small as he likes. You answer with an $N$, knowing his $\eps$. He picks any $n\ge N$. You win if
$|a_n-a|<\eps$. A proof is a strategy that wins every time: a recipe $\eps\mapsto N$. That is the whole reason $N$ may depend
on $\eps$ and never on $n$. You move before he picks $n$.
\aside{Weierstraß turned limits into this game in his Berlin lectures around 1860. Nearly two centuries of quarrelling about
infinitesimals ended there.}

To \emph{refute} a statement, switch sides. The negation flips every quantifier and negates the last line. An implication
$A\Rightarrow B$ fails exactly when $A$ holds and $B$ does not. Every ``zeigen Sie, dass \dots\ nicht'' in the test is one of
these rows.

\begin{panel}{Table}{definitions and their negations}{Verneinung}
\begin{negtable}
\negrow{Convergence $a_n\to a$ \de{(Konvergenz)}}
 {\forall\eps>0\ \exists N\ \forall n\ge N:\ |a_n-a|<\eps}
 {\exists\eps>0\ \forall N\ \exists n\ge N:\ |a_n-a|\ge\eps}
\negrow{Cauchy sequence}
 {\forall\eps>0\ \exists N\ \forall m,n\ge N:\ |a_m-a_n|<\eps}
 {\exists\eps>0\ \forall N\ \exists m,n\ge N:\ |a_m-a_n|\ge\eps}
\negrow{$M$ bounded above \de{(nach oben beschränkt)}}
 {\exists K\ \forall x\in M:\ x\le K}
 {\forall K\ \exists x\in M:\ x>K}
\negrow{An upper bound $s$ of $M$ is the supremum}
 {\forall\eps>0\ \exists x\in M:\ x>s-\eps}
 {\exists\eps>0\ \forall x\in M:\ x\le s-\eps}
\negrow{$f$ continuous at $x_0$ \de{(stetig in $x_0$)}}
 {\forall\eps>0\ \exists\delta>0\ \forall x\in D:\ |x-x_0|<\delta\Rightarrow|f(x)-f(x_0)|<\eps}
 {\exists\eps>0\ \forall\delta>0\ \exists x\in D:\ |x-x_0|<\delta\ \text{and}\ |f(x)-f(x_0)|\ge\eps}
\negrow{$f$ uniformly continuous on $D$}
 {\forall\eps>0\ \exists\delta>0\ \forall x,y\in D:\ |x-y|<\delta\Rightarrow|f(x)-f(y)|<\eps}
 {\exists\eps>0\ \forall\delta>0\ \exists x,y\in D:\ |x-y|<\delta\ \text{and}\ |f(x)-f(y)|\ge\eps}
\negrow{$f$ injective}
 {\forall x_1,x_2:\ f(x_1)=f(x_2)\Rightarrow x_1=x_2}
 {\exists x_1\ne x_2:\ f(x_1)=f(x_2)}
\negrow{$f\colon M\to N$ surjective}
 {\forall y\in N\ \exists x\in M:\ f(x)=y}
 {\exists y\in N\ \forall x\in M:\ f(x)\ne y}
\end{negtable}
\end{panel}

\why{Why does the order $\forall\eps\ \exists N$ matter so much?}{Swap it to $\exists N\ \forall\eps$ and you must name one
$N$ that works for every $\eps$ at once. Then $|a_n-a|<\eps$ for all $\eps>0$, so $a_n=a$ for every $n\ge N$: the sequence is
eventually constant. Same symbols, a different statement. The order is the content.}

\subsection{Every proof is written twice}
Here is the difficulty you described, in one sentence: the proof must \emph{choose} $N$ in its second line, but you only learn
which $N$ works at the end of the computation. So you compute twice. The scratch work runs backwards, from the goal to a
condition on $n$. The proof runs forwards and looks like magic, because the magic happened on the scratch paper. In the exam you
may show the scratch as a \emph{Vorüberlegung}; the proof must still stand on its own.

\par\needspace{78mm}
\begin{center}
\begin{tikzpicture}
\begin{axis}[klplot,width=.88\linewidth,height=52mm,xmin=0,xmax=41,ymin=0.6,ymax=2.45,
  xlabel={$n$},ylabel={$a_n$},xtick={1,10,21,30,40},ytick={1,1.7,2,2.3},
  yticklabels={$1$,$2-\varepsilon$,$2$,$2+\varepsilon$},ylabel style={rotate=-90}]
  \fill[pattern={dither25},pattern color=klink] (axis cs:21,1.7) rectangle (axis cs:41,2.3);
  \draw[klink,line width=.5pt] (axis cs:21,1.7)--(axis cs:41,1.7) (axis cs:21,2.3)--(axis cs:41,2.3);
  \draw[kldim,dashed,line width=.5pt] (axis cs:0,2)--(axis cs:41,2);
  \draw[kldim,dashed,line width=.5pt] (axis cs:21,0.6)--(axis cs:21,2.45) node[above,text=klink]{$N=21$};
  \addplot[only marks,mark=square*,mark size=1.4pt,color=klhi,samples at={21,...,40}] {(2*x^2+1)/(x^2+3*x)};
  \addplot[only marks,mark=square,mark size=1.4pt,color=kldim,samples at={1,...,20}] {(2*x^2+1)/(x^2+3*x)};
\end{axis}
\end{tikzpicture}
\end{center}
\plotcaption{The example below, $a_n=\frac{2n^2+1}{n^2+3n}\to2$, with $\eps=0.3$. The proof's $N=21$ works: every later
term sits in the band. The sequence actually enters the band
at $n=17$. Nobody cares. The definition asks for \emph{some} $N$, so crude estimates are allowed, and they are faster.}

\begin{panel}[unbreakable]{Example}{two documents}{}
Claim: $a_n=\dfrac{2n^2+1}{n^2+3n}\to 2$.
\begin{twodocs}
Goal: $|a_n-2|<\eps$. Make one fraction:
\[\begin{aligned}
|a_n-2|&=\left|\frac{2n^2+1-2n^2-6n}{n^2+3n}\right|\\
&=\frac{6n-1}{n^2+3n}.
\end{aligned}\]
Bigger and simpler:
\[ \frac{6n-1}{n^2+3n}<\frac{6n}{n^2}=\frac6n . \]
Enough: $\frac6n<\eps$, that is, $n>\frac6\eps$.
&
Let $\eps>0$. By the Archimedean property there is $N\in\N$ with $N>\frac6\eps$. For every $n\ge N$,
\[\begin{aligned}
|a_n-2|&=\frac{6n-1}{n^2+3n}\\
&<\frac6n\le\frac6N<\eps .
\end{aligned}\]
Hence $a_n\to2$.\hfill\qedsymbol
\end{twodocs}
\end{panel}


\subsection{The toolbox}
Nearly every estimate in chapters 2--4 is one of these. When you are stuck, go down the list.

\medskip
\begin{twocol}{66mm}{Tool}{Reach for it when}
\row{$|x+y|\le|x|+|y|$ \ and \ $\bigl||x|-|y|\bigr|\le|x-y|$}{a difference must be split; $|a_n|\to|a|$}
\row{numerator up, denominator down}{$\frac{6n-1}{n^2+3n}<\frac{6n}{n^2}$: simpler and bigger}
\row{$n^2-5\ge\frac{n^2}{2}$ for $n\ge4$}{a minus sign in a denominator: keep half of the top power}
\row{$\sqrt a-\sqrt b=\dfrac{a-b}{\sqrt a+\sqrt b}$}{a difference of roots}
\row{$x^n-y^n=(x-y)(x^{n-1}+x^{n-2}y+\dots+y^{n-1})$}{powers; $\frac{x^N-1}{x-1}$; monotonicity of $x^n$}
\row{$(1+h)^n\ge1+nh$ \ for $h\ge-1$}{$q^n\to0$ for $|q|<1$; $(1+\frac1n)^n\ge2$}
\row{$(1+h)^n\ge\binom n2h^2$ \ for $h\ge0$, $n\ge2$}{$\sqrt[n]n\to1$}
\row{$a_nb_n-ab=(a_n-a)b_n+a(b_n-b)$}{products; then $\frac\eps2+\frac\eps2$}
\row{$a_n\to a>0\ \Rightarrow\ a_n>\frac a2$ for large $n$}{quotients: play $\eps=\frac a2$}
\row{restrict $\delta\le1$}{$\eps$-$\delta$: it freezes the other factor, e.g.\ $|x+x_0|<1+2|x_0|$}
\row{$\frac1{k^2}\le\frac1{k(k-1)}=\frac1{k-1}-\frac1k$}{bounding $\sum\frac1{k^2}$ by a telescope}
\row{$(1+\frac1n)^n\le3$}{the Skript has no limit $e$; this bound replaces it}
\row{$\exp x\ge1+x+\frac{x^2}2$ for $x\ge0$}{$\frac{x}{\exp x}\to0$: exp beats every power}
\end{twocol}

\medskip
Four of these need a line of proof, because the test will not let you assume them.

\emph{Reverse triangle inequality.} $|x|=|(x-y)+y|\le|x-y|+|y|$, so $|x|-|y|\le|x-y|$. Swap $x$ and $y$.

\emph{Keep half.} $n^2-5\ge\frac{n^2}2\iff n^2\ge10$, which holds for $n\ge4$.

\emph{Binomial lower bound.} In $(1+h)^n=\sum_k\binom nk h^k$ every term is $\ge0$. Keep only $k=2$.

\emph{The bound $(1+\frac1n)^n\le3$.} By the binomial theorem $(1+\frac1n)^n=\sum_{k=0}^n\binom nk\frac1{n^k}$, and for $k\ge1$
\[
  \binom nk\frac1{n^k}=\frac{n(n-1)\cdots(n-k+1)}{k!\;n^k}\le\frac1{k!}\le\frac1{2^{k-1}} .
\]
So the sum is at most $1+\bigl(1+\frac12+\frac14+\cdots\bigr)=3$.

% =====================================================================================
\section{The real numbers}

\opening{T}{he Skript builds} $\R$ in three layers. The field axioms K1--K9 give you arithmetic. The order axioms O1--O4 give
you inequalities. $\Q$ satisfies both layers and is useless for analysis, because it is full of holes. The third layer is one
sentence, and it fills them.
\aside{Dedekind invented the cuts of your Skript's Exkurs in 1858, teaching calculus in Zürich, because he refused to prove that
bounded increasing sequences converge by pointing at a drawing.}

\begin{keystone}{Completeness}{Vollständigkeit}
\textbf{The supremum axiom.} Every nonempty subset of $\R$ that is bounded above has a least upper bound, its supremum.

\smallskip\noindent
Everything hard in chapters 1--4 stands on this sentence, and every item below is \emph{false} in $\Q$: $\sqrt2$ exists; nested
intervals share a point; monotone bounded sequences converge; bounded sequences have convergent subsequences; Cauchy sequences
converge; a continuous $f$ on $[a,b]$ takes every value between $f(a)$ and $f(b)$, attains a maximum and a minimum, and is
uniformly continuous.
\end{keystone}

\par\needspace{100mm}
\begin{center}
\begin{tikzpicture}[x=1mm,y=1mm,
  box/.style={draw=kllow,leds,align=center,font=\fell\fontsize{9}{10.5}\selectfont,text=klink,
    minimum width=27mm,minimum height=8.5mm,inner sep=2pt},
  top/.style={box,text=klhi,draw=kldim},
  arr/.style={-{Stealth[length=4pt]},kldim,line width=.5pt}]
  \node[top] (S) at (64,0) {Supremum axiom};
  \node[box] (A) at (15,-16) {Archimedes};
  \node[box] (I) at (48,-16) {Nested intervals};
  \node[box] (M) at (81,-16) {Monotone\\convergence};
  \node[box] (R) at (114,-16) {$\sqrt a$ exists};
  \node[box] (D) at (15,-32) {$\mathbb Q$ is dense};
  \node[box] (B) at (48,-32) {Bolzano--\\Weierstraß};
  \node[box] (L) at (81,-32) {Leibniz};
  \node[box] (Z) at (114,-32) {Intermediate\\value theorem};
  \node[box] (C) at (15,-48) {Cauchy\\criterion};
  \node[box] (E) at (48,-48) {Extreme\\value theorem};
  \node[box] (H) at (81,-48) {Heine--Cantor};
  \node[box] (U) at (114,-48) {Inverse\\function theorem};
  \node[box] (T) at (15,-64) {Absolute conv.,\\majorant, ratio};
  \node[box] (F) at (48,-64) {Continuous\\extension};
  \node[box] (P) at (81,-64) {Step-function\\approximation};
  \node[box] (X) at (114,-64) {$\exp$ and $\log$\\bijective};
  \foreach \a/\b in {S/A,S/I,S/M,S/R,A/D,I/B,M/L,M/Z,B/C,B/E,B/H,Z/U,C/T,C/F,H/P,E/P,U/X}
    \draw[arr] (\a)--(\b);
\end{tikzpicture}
\end{center}
\plotcaption{Each arrow reads ``is proved with''. Asked to prove a theorem, first ask which arrow points into it. That is the
idea, and the idea is most of the marks.}

\why{Why is the Archimedean property not on the list of things that fail in $\Q$?}{Because it holds in $\Q$: for
$r=\frac pm>0$ with $p,m\in\N$, the integer $p+1$ exceeds $r$. The Skript derives it from the supremum axiom only because a
real number is not known to be a fraction. What $\Q$ lacks is completeness: $\{q\in\Q: q^2<2\}$ is bounded and has no least
upper bound \emph{in} $\Q$, and every theorem on the list falls into that hole.}

\subsection{The terms of chapter 1}
The cheapest points on the test. For each: the definition, an example, a non-example.

\entry{Sets}{Mengen}
$M\subseteq N$ means $x\in M\Rightarrow x\in N$. $M=N$ means both inclusions, so equality is proved as two inclusions, or as a
chain of $\iff$ on membership. De Morgan: $(M\cup N)^c=M^c\cap N^c$ and $(M\cap N)^c=M^c\cup N^c$.

\entry{Map}{Abbildung}
assigns to each $x\in M$ exactly one $f(x)\in N$. Image $f(A)=\{f(x):x\in A\}$. Preimage $f^{-1}(B)=\{x\in M: f(x)\in B\}$;
it exists for every map, inverse or not.
\eg $f(x)=x^2$ on $\R$: $f([-1,2])=[0,4]$ and $f^{-1}([1,4])=[-2,-1]\cup[1,2]$.

\entry{Injective, surjective, bijective}{injektiv, surjektiv, bijektiv}
$f(x_1)=f(x_2)\Rightarrow x_1=x_2$; every $y\in N$ is hit; both. $f$ is bijective iff some $g$ has $g\circ f=\mathrm{id}_M$
and $f\circ g=\mathrm{id}_N$. Then $g=f^{-1}$, and $(g\circ f)^{-1}=f^{-1}\circ g^{-1}$.
\eg $x\mapsto x^2$ from $[0,\infty)$ to $[0,\infty)$. \nonex the same rule from $\R$ to $\R$: neither.

\entry{Field}{Körper}
a set with $+$ and $\cdot$ and elements $0\ne1$ satisfying K1--K9: associativity, commutativity, neutral elements, inverses,
distributivity. Inverses are then unique, and $0\cdot x=0$ follows from $0x=(0+0)x=0x+0x$.
\eg $\Q$, $\R$, $\mathbb C$, $\mathbb Z_2=\{0,1\}$. \nonex $\mathbb Z$, where $2$ has no inverse.

\entry{Ordered field}{angeordneter Körper}
a field with $<$ satisfying O1--O4: trichotomy, transitivity, and $<$ survives adding $z$ and multiplying by $z>0$. Then
$x^2>0$ for $x\ne0$, in particular $1>0$. Multiplying by $z<0$ flips an inequality; $0\le x<y$ gives $x^2<y^2$; and
$|x|\le K\iff-K\le x\le K$. \nonex $\mathbb C$: in any order $i^2=-1$ would be positive.

\entry{Absolute value}{Betrag}
$|x|=x$ for $x\ge0$ and $-x$ otherwise. $|xy|=|x||y|$, $|x+y|\le|x|+|y|$, and the workhorse
\[ |x-a|<r\iff a-r<x<a+r . \]
\eg $|2x-1|<3\iff-1<x<2$.

\entry{Interval}{Intervall}
a set $I$ with: $x,y\in I$ and $x\le z\le y$ imply $z\in I$.
\eg $[0,1]$, $(0,\infty)$, $\{a\}$. \nonex $\{0,1\}$ and $[0,1]\cup[2,3]$.

\entry{Bounds, maximum, minimum}{Schranken, Maximum, Minimum}
$K$ is an upper bound of $M$ if $x\le K$ for all $x\in M$. A maximum is an upper bound that lies \emph{in} $M$. $M$ is bounded
iff $|x|\le K$ for all $x\in M$, for some $K$.
\eg $\max\{1,3,5\}=5$. \nonex $(0,1)$ has no maximum.

\entry{Supremum, infimum}{Supremum, Infimum}
$\sup M$ is the least upper bound: an upper bound $s$ such that
\[ \forall\eps>0\ \ \exists x\in M:\ x>s-\eps . \]
$\inf M=-\sup(-M)$. If $\max M$ exists, it is $\sup M$.
\eg $\sup(0,1)=1\notin(0,1)$. \nonex $\N$ has no supremum in $\R$.

\entry{Archimedean property}{Archimedische Eigenschaft}
for every $r\in\R$ some $n\in\N$ has $n>r$; equivalently, for every $\eps>0$ some $\frac1n<\eps$. Every $N$ in Part~IV is
produced by this sentence.

\entry{Countable}{abzählbar}
$M$ is countably infinite if a bijection $\N\to M$ exists; ``countable'' often also admits finite sets.
\eg $\mathbb Z$, $\Q$, $\N\times\N$. \nonex $\R$ and $[0,1]$ are uncountable \de{(überabzählbar)}.

\entry{Nested intervals}{Intervallschachtelung}
closed bounded intervals $I_n=[a_n,b_n]$ with $I_{n+1}\subseteq I_n$. Their intersection is never empty, and it is a single point
when $b_n-a_n\to0$. \nonex $(0,\frac1n)$: open, and the intersection is empty.

\entry{Complex numbers}{komplexe Zahlen}
$z=x+iy$ with $i^2=-1$, $\bar z=x-iy$, $|z|=\sqrt{x^2+y^2}$, $z\bar z=|z|^2$. To divide, multiply by the conjugate:
$\frac1{1-2i}=\frac{1+2i}{5}$.

% ------------------------------------------------------------------------------------
\subsection{Move 01\quad Induction}
\when{``Zeigen Sie für alle $n\in\N$ (bzw.\ $n\ge n_0$)'': a sum formula, an inequality, a divisibility, anything indexed by $n$.}

\begin{panel}[unbreakable]{Skeleton}{induction}{vollständige Induktion}
\begin{enumerate}
  \item State $A(n)$ exactly, with its starting index $n_0$.
  \item \textbf{Base} \de{(Induktionsanfang)}: check $A(n_0)$ by computing both sides.
  \item \textbf{Hypothesis} \de{(Induktionsvoraussetzung)}: $A(n)$ holds for one fixed $n\ge n_0$.
  \item \textbf{Step} \de{(Induktionsschritt)}: start from one side of $A(n+1)$, find $A(n)$ inside it, use the hypothesis,
        walk to the other side. For an inequality, replace the $A(n)$-part by its bound and prove the small leftover
        inequality separately.
  \item Close: ``Nach dem Induktionsprinzip gilt $A(n)$ für alle $n\ge n_0$.''
\end{enumerate}
\end{panel}

\begin{example}{an inequality by induction}
Claim: $2^n\ge n^2$ for all $n\ge4$.

\emph{Base.} $2^4=16=4^2$.

\emph{Step.} Let $n\ge4$ with $2^n\ge n^2$. Then $2^{n+1}=2\cdot2^n\ge2n^2$. Leftover: is $2n^2\ge(n+1)^2$? That is
$n^2\ge2n+1$, that is $(n-1)^2\ge2$, true because $n\ge4$ gives $(n-1)^2\ge9$. So
\[ 2^{n+1}\ge2n^2\ge(n+1)^2 . \]
The claim fails for $n=3$ ($8<9$). The base is not decoration: it is where $n_0$ comes from.
\end{example}

\trap{Opening the step with the goal $2^{n+1}\ge(n+1)^2$ and manipulating it until something true appears. That proves the
converse. Start on one side and walk to the other. (Your UE01/8.)}
\trap{Using the step where it fails. Say where $n\ge n_0$ enters; here the leftover needed $n\ge3$.}

\begin{yourturn}
\task{t:fac}$\sum_{k=1}^n k\cdot k!=(n+1)!-1$ for all $n\in\N$.
\task{t:fac2}$n!>2^n$ for all $n\ge4$.
\par\smallskip\noindent Then Vorbereitung~§1 with the sheet closed.
\end{yourturn}

% ------------------------------------------------------------------------------------
\subsection{Move 02\quad Supremum and infimum}
\when{``Bestimmen Sie Supremum, Infimum, Maximum und Minimum (falls existent)''; ``Zeigen Sie $\sup(A+B)=\sup A+\sup B$''.}

\begin{panel}[unbreakable]{Skeleton}{supremum}{Supremum nachweisen}
\begin{enumerate}
  \item \textbf{Guess.} Compute the first elements, split by parity, draw the number line.
  \item \textbf{Upper bound.} Show $x\le s$ for every $x\in M$.
  \item \textbf{Least.} For every $\eps>0$ produce an $x\in M$ with $x>s-\eps$. The Archimedean property picks the index.
  \item \textbf{Attained?} If $s\in M$, then $s=\max M$. If not, say so and why (a strict inequality for every element).
  \item \textbf{Infimum:} the mirror image, $x<i+\eps$.
\end{enumerate}
\end{panel}

\begin{example}{a set with neither maximum nor minimum}
$M=\bigl\{(-1)^n\frac{n}{n+1}: n\in\N\bigr\}=\{-\frac12,\ \frac23,\ -\frac34,\ \frac45,\dots\}$.

\emph{Guess.} Even indices climb towards $1$, odd ones fall towards $-1$.

\emph{Bounds.} $\bigl|(-1)^n\frac n{n+1}\bigr|=\frac n{n+1}<1$, so every element lies strictly between $-1$ and $1$. So $1$ is
an upper bound and $-1$ a lower bound, and neither is attained: no maximum, no minimum.

\emph{Least.} Let $\eps>0$. Choose $N\in\N$ with $N>\frac1\eps$ and take the even index $n=2N$. Then
\[ \frac{n}{n+1}=1-\frac1{n+1}>1-\frac1N>1-\eps . \]
So $\sup M=1$. The odd index $n=2N+1$ gives $-\frac n{n+1}<-1+\eps$ the same way, so $\inf M=-1$.
\end{example}

\par\needspace{40mm}
\begin{center}
\begin{tikzpicture}[x=50mm,y=1mm]
  \fill[pattern={dither25},pattern color=klink] (0.75,-2.2) rectangle (1,2.2);
  \draw[kldim,line width=.6pt] (-1.18,0)--(1.18,0);
  \foreach \n in {1,...,16}{\pgfmathsetmacro\v{(-1)^\n*\n/(\n+1)}\fill[klhi] (\v,0) circle (1.3pt);}
  \draw[klink] (1,-3)--(1,3) node[above]{$\sup M=1$};
  \draw[klink] (-1,-3)--(-1,3) node[above]{$\inf M=-1$};
  \draw[kldim,dashed] (0.75,-3)--(0.75,3) node[above]{$1-\eps$};
  \draw[kldim] (0,-1.5)--(0,1.5) node[below=4mm]{$0$};
\end{tikzpicture}
\end{center}
\plotcaption{For $\eps=\frac14$ the window $(1-\eps,1]$ contains elements of $M$, here every even-indexed one from $n=4$ on.
That is what step 3 proves for every $\eps$. The window never contains $1$ itself.}

\begin{example}{supremum of a sum}
Let $A,B\subseteq\R$ be nonempty and bounded above, $A+B=\{a+b: a\in A,\ b\in B\}$. Then $\sup(A+B)=\sup A+\sup B$.

\emph{Upper bound.} $a+b\le\sup A+\sup B$ for all $a,b$.

\emph{Least.} Let $\eps>0$. Pick $a\in A$ with $a>\sup A-\frac\eps2$ and $b\in B$ with $b>\sup B-\frac\eps2$. Then
$a+b>\sup A+\sup B-\eps$.\hfill\qedsymbol

\noindent The $\frac\eps2$-split appears here for the first time. It will not be the last.
\end{example}

\trap{Stopping after the upper bound. ``$1$ is an upper bound and the terms get close'' describes; step 3 proves, and step 3 is
where the marks are.}
\trap{Confusing supremum and maximum. A bounded nonempty set always has a supremum, a maximum only if the supremum belongs to
it.}

\begin{yourturn}
\task{t:supmn}$M=\{\frac1m+\frac1n: m,n\in\N\}$. Find $\sup M$, $\inf M$, $\max M$, $\min M$ and prove each.
\task{t:supc}$A\ne\emptyset$ bounded above and $c>0$. Show $\sup(cA)=c\sup A$, where $cA=\{ca: a\in A\}$.
\par\smallskip\noindent Then Vorbereitung~§2 cold, and your Übung~2, Aufgabe~2 cold.
\end{yourturn}

% ------------------------------------------------------------------------------------
\subsection{Move 03\quad Bijectivity and the inverse}
\when{``Untersuchen Sie $f$ auf Injektivität, Surjektivität und Bijektivität und bestimmen Sie gegebenenfalls $f^{-1}$.''}

\begin{panel}[unbreakable]{Skeleton}{bijectivity}{Bijektivität}
\begin{enumerate}
  \item Write down the domain and the codomain first. The same formula between other sets is another function.
  \item \textbf{Injective:} assume $f(x_1)=f(x_2)$ and derive $x_1=x_2$. Or show $f$ strictly monotone.
  \item \textbf{Surjective:} take an arbitrary $y$ in the codomain, solve $f(x)=y$, then check that this $x$ lies in the domain.
  \item \textbf{Inverse:} the $x$ from step 3 is $f^{-1}(y)$. Check $f(f^{-1}(y))=y$.
  \item \textbf{To refute:} two points with the same value; one $y$ that nobody hits.
\end{enumerate}
\end{panel}

\begin{example}{a rational function between punctured lines}
$f\colon\R\setminus\{-1\}\to\R\setminus\{1\}$, $f(x)=\dfrac{x-1}{x+1}$.

First, the codomain is honest: $f(x)=1-\frac2{x+1}$, and $\frac2{x+1}\ne0$, so $f(x)\ne1$.

\emph{Injective.} $f(x_1)=f(x_2)\Rightarrow\frac2{x_1+1}=\frac2{x_2+1}\Rightarrow x_1=x_2$.

\emph{Surjective.} Let $y\ne1$. For $x\ne-1$:
\[\begin{aligned}
\frac{x-1}{x+1}=y&\iff x-1=yx+y\\
&\iff x(1-y)=1+y\\
&\iff x=\frac{1+y}{1-y},
\end{aligned}\]
allowed because $y\ne1$. And this $x$ is not $-1$: $\frac{1+y}{1-y}=-1$ would mean $1+y=y-1$.

\emph{Inverse.} $f^{-1}\colon\R\setminus\{1\}\to\R\setminus\{-1\}$, $f^{-1}(y)=\dfrac{1+y}{1-y}$.
\end{example}

\trap{Solving $f(x)=y$ without checking the domain. List every condition the formula needs (no zero denominator, nothing
negative under a root, the interval) and translate each into a condition on $y$. (Your UE01/5c.)}
\trap{Reading $f^{-1}(B)$ as an inverse map. The preimage exists for every $f$; the inverse only for bijections.}

\begin{yourturn}
\task{t:bij}$g\colon[0,\infty)\to[1,\infty)$, $g(x)=x^2+2x+1$. Bijective? If so, find $g^{-1}$.
\task{t:comp}$f\colon M\to N$, $g\colon N\to P$. Show: $g\circ f$ injective $\Rightarrow f$ injective, and $g\circ f$
surjective $\Rightarrow g$ surjective. Find an example with $g\circ f$ bijective but $g$ not injective.
\end{yourturn}

% =====================================================================================
\section{Sequences}

\opening{A}{sequence} is a map $a\colon\N\to\R$, and convergence is the game of Part~II. Chapter~2 adds three ways to win it
without knowing the limit in advance: monotonicity, subsequences, and the Cauchy property. All three are the supremum axiom
again.

\subsection{The terms of chapter 2}

\entry{Sequence}{Folge}
a map $\N\to\R$, $n\mapsto a_n$, possibly defined recursively. \eg $a_1=1$, $a_{n+1}=\sqrt{1+a_n}$.

\entry{Subsequence}{Teilfolge}
$(a_{n_k})_k$ with $n_1<n_2<\cdots$. Then $n_k\ge k$, the fact behind every subsequence proof.
\eg $(a_{2k})$. \nonex $(a_1,a_1,a_2,a_2,\dots)$: the indices do not strictly increase.

\entry{Convergence, null sequence}{Konvergenz, Nullfolge}
$a_n\to a$ iff $\forall\eps>0\ \exists N\ \forall n\ge N: |a_n-a|<\eps$; a null sequence converges to $0$. $a_n\to a$ iff
$(a_n-a)$ is a null sequence. The limit is unique. \nonex $(-1)^n$.

\entry{Accumulation point of a sequence}{Häufungspunkt einer Folge}
the limit of some subsequence. A convergent sequence has exactly one, its limit.
\eg $(-1)^n$ has $\pm1$. Not to be confused with an accumulation point of a \emph{set} $D$ (Part~VI).

\entry{Bounded, monotone}{beschränkt, monoton}
$|a_n|\le K$ for all $n$; $a_n\le a_{n+1}$ for all $n$ (increasing). Convergent $\Rightarrow$ bounded, never the reverse.

\entry{Cauchy sequence}{Cauchy-Folge}
$\forall\eps>0\ \exists N\ \forall m,n\ge N: |a_m-a_n|<\eps$. In $\R$: Cauchy $\iff$ convergent.
\eg $\frac1n$. \nonex $\sqrt n$, although $\sqrt{n+1}-\sqrt n\to0$.

\entry{Definite divergence}{bestimmte Divergenz}
$a_n\to+\infty$ iff $\forall K\ \exists N\ \forall n\ge N: a_n>K$. This is a kind of divergence, not convergence: $\infty\notin\R$.

\begin{panel}{Table}{limits you may quote, with their one-line reason}{Standardgrenzwerte}
\begin{twocol}{44mm}{Limit}{Reason}
\row{$\frac cn\to0$}{Archimedes: $n>\frac c\eps$}
\row{$q^n\to0$ \ for $|q|<1$}{write $\frac1{|q|}=1+h$; Bernoulli gives $|q|^n\le\frac1{1+nh}$}
\row{$\sqrt[n]a\to1$ \ for $a>0$}{Skript: $\sqrt[n]a=\exp\bigl(\frac{\log a}{n}\bigr)\to\exp0$}
\row{$\sqrt[n]n\to1$}{$h_n=\sqrt[n]n-1$: $n\ge\binom n2h_n^2$, so $0\le h_n\le\sqrt{\frac2{n-1}}$}
\row{$n^kq^n\to0$, \ $\frac{c^n}{n!}\to0$}{their series converge by the ratio test, so the terms tend to $0$}
\row{$\sqrt{n+1}-\sqrt n\to0$}{conjugate: $=\frac1{\sqrt{n+1}+\sqrt n}<\frac1{2\sqrt n}$}
\end{twocol}
\end{panel}

% ------------------------------------------------------------------------------------
\subsection{\texorpdfstring{Move 04\quad Convergence with $\eps$ and $N$}{Move 04: convergence with epsilon and N}}
\when{``Zeigen Sie mit der Definition (dem $\eps$-Kriterium), dass $a_n\to a$.'' If the question says only ``Bestimmen Sie
den Grenzwert'', the limit laws are allowed: divide by the highest power and cite the Rechenregeln.}

\begin{panel}[unbreakable]{Skeleton}{epsilon-N}{Konvergenz nach Definition}
\begin{twodocs}
\begin{enumerate}
  \item Write $|a_n-a|$ as one fraction; simplify.
  \item Estimate upwards to $\frac Cn$ (or $\frac C{\sqrt n}$, or $Cq^n$). Record every condition you impose, such as $n\ge3$.
  \item Solve $\frac Cn<\eps$ for $n$.
\end{enumerate}
&
``Sei $\eps>0$. Wähle $N\in\N$ mit $N\ge n_0$ und $N>\frac C\eps$ (Archimedes). Für alle $n\ge N$ gilt
\[ |a_n-a|=\dots\le\frac Cn\le\frac CN<\eps .\text{''} \]
\end{twodocs}
\end{panel}

\begin{example}{a minus sign in the denominator}
Claim: $a_n=\dfrac{n^2+2n}{2n^2-5}\to\dfrac12$.

\emph{Scratch.} For $n\ge2$ the denominator is positive, and
\[ \Bigl|a_n-\frac12\Bigr|=\frac{|2n^2+4n-2n^2+5|}{2(2n^2-5)}=\frac{4n+5}{4n^2-10}. \]
Numerator up: $4n+5\le6n$ for $n\ge3$. Denominator down, keeping half: $4n^2-10\ge2n^2$ for $n\ge3$. So for $n\ge3$
\[ \Bigl|a_n-\frac12\Bigr|\le\frac{6n}{2n^2}=\frac3n . \]
\emph{Proof.} Let $\eps>0$. Choose $N\in\N$ with $N\ge3$ and $N>\frac3\eps$. For every $n\ge N$:
\[ \Bigl|a_n-\frac12\Bigr|=\frac{4n+5}{4n^2-10}\le\frac3n\le\frac3N<\eps .\]
\end{example}

\noindent Roots go the same way after the conjugate. For $b_n=\sqrt{n+1}-\sqrt n$:
\[ |b_n-0|=\frac1{\sqrt{n+1}+\sqrt n}<\frac1{2\sqrt n}<\eps\quad\text{as soon as}\quad n>\frac1{4\eps^2}. \]

\trap{$N$ depending on $n$ (``wähle $N=n+1$''). $N$ is chosen before any $n$ exists.}
\trap{A chain that changes direction. Every link is $=$, $\le$ or $<$, all pointing the same way, and the chain ends in $<\eps$.}
\trap{Forgetting your side conditions. You used $n\ge3$, so $N\ge3$ as well.}

\begin{yourturn}
\task{t:epsA}With the definition: $\dfrac{3n-2}{n+4}\to3$.
\task{t:epsB}With the definition: $\dfrac{n+(-1)^n}{n^2}\to0$.
\task{t:abs}$a_n\to a\Rightarrow|a_n|\to|a|$. Does the converse hold?
\task{t:sqrt}$a_n\ge0$ and $a_n\to a$ $\Rightarrow$ $\sqrt{a_n}\to\sqrt a$ (your Übung~3/5, cold).
\end{yourturn}

% ------------------------------------------------------------------------------------
\subsection{Move 05\quad Subsequences and accumulation points}
\when{``Bestimmen Sie alle Häufungspunkte''; ``Ist die Folge konvergent?'' for anything that oscillates: $(-1)^n$, case
distinctions, periodic patterns.}

\begin{panel}[unbreakable]{Skeleton}{accumulation points}{Häufungspunkte}
\begin{enumerate}
  \item Find the period $p$ of the pattern (parity: $p=2$).
  \item For each remainder class $r$, compute $\lim_k a_{pk+r}$. Each such limit is an accumulation point.
  \item No others: a convergent subsequence has infinitely many indices in some class; the part of it in that class tends to
        that class's limit, and a subsequence of a convergent sequence keeps the limit.
  \item Two different accumulation points $\Rightarrow$ divergent.
\end{enumerate}
\end{panel}

\begin{example}{a sign pattern of period four}
$a_n=(-1)^{n(n+1)/2}\cdot\dfrac n{n+1}$. The signs for $n=1,2,3,\dots$ run $-,-,+,+,-,-,+,+,\dots$ Write $n=4k+r$ and look at
the parity of $\frac{n(n+1)}2$:
\[
\begin{array}{c|c|c}
n & \frac{n(n+1)}2 & \text{sign}\\\hline
4k+1 & (4k+1)(2k+1) & \text{odd: }-\\
4k+2 & (2k+1)(4k+3) & \text{odd: }-\\
4k+3 & (4k+3)(2k+2) & \text{even: }+\\
4k+4 & (2k+2)(4k+5) & \text{even: }+
\end{array}
\]
Since $\frac n{n+1}\to1$, the four classes tend to $-1,-1,1,1$. The accumulation points are exactly $\pm1$, and the sequence
diverges.
\end{example}

\trap{``Only one accumulation point, so it converges.'' False when the sequence is unbounded: $a_n=n^{(-1)^n}$ is $\frac1n$
for odd $n$ and $n$ for even $n$. Its only accumulation point is $0$. For \emph{bounded} sequences the statement is true
(Part~VII).}

\begin{yourturn}
\task{t:hp}$a_n=(-1)^n\dfrac n{n+1}+\dfrac{1+(-1)^n}2$. All accumulation points? Convergent?
\par\smallskip\noindent Then Vorbereitung~§5 cold.
\end{yourturn}

% ------------------------------------------------------------------------------------
\subsection{Move 06\quad Monotone convergence and Cauchy}
\when{a recursive sequence $a_{n+1}=g(a_n)$; ``Zeigen Sie, dass $(a_n)$ konvergiert'' with no limit in sight; partial sums of
positive terms; ``Zeigen Sie direkt, dass $(a_n)$ eine Cauchy-Folge ist''.}

\begin{panel}[unbreakable]{Skeleton}{recursive sequences}{rekursive Folgen}
\begin{enumerate}
  \item \textbf{Guess} $L$ from $L=g(L)$; compute three terms to see which root, and whether the sequence climbs or falls.
  \item \textbf{Bound} by induction, for example $0<a_n<L$.
  \item \textbf{Monotone:} factor $a_{n+1}-a_n$ and read off its sign from step 2, or induct.
  \item Monotone and bounded $\Rightarrow$ convergent \de{(Monotoniekriterium)}. Call the limit $L$.
  \item Only now: $a_{n+1}\to L$ as well, so $L=g(L)$ by the limit laws. The bound from step 2 picks the root.
\end{enumerate}
\end{panel}

\begin{example}{a recursion with two candidate limits}
$a_1=\frac12$, $a_{n+1}=\dfrac{a_n^2+2}3$.

\emph{Guess.} $L=\frac{L^2+2}3\iff L^2-3L+2=0\iff L\in\{1,2\}$. Terms: $0.5,\ 0.75,\ 0.85,\ 0.91$. Climbing towards $1$.

\emph{Bound.} Claim $0<a_n<1$. True for $n=1$. If $0<a_n<1$, then $a_n^2<1$, so $0<a_{n+1}<\frac{1+2}3=1$.

\emph{Monotone.}
\[ a_{n+1}-a_n=\frac{a_n^2-3a_n+2}3=\frac{(a_n-1)(a_n-2)}3>0, \]
because both factors are negative by the bound.

\emph{Limit.} Increasing and bounded above, so $a_n\to L$. From $a_{n+1}=\frac{a_n^2+2}3$ in the limit, $L\in\{1,2\}$. Since
$a_n<1$ for all $n$, $L\le1$. So $L=1$.
\end{example}

\why{Why is step 5 alone not a proof?}{It assumes the limit exists. For $a_1=1$, $a_{n+1}=2a_n$ the same computation gives
$L=2L$, so $L=0$, while $a_n=2^{n-1}\to\infty$. The equation $L=g(L)$ says what the limit must be \emph{if} there is one. Steps
2 to 4 prove that there is.}

\noindent\textbf{Cauchy, directly.} Bound $|a_m-a_n|$ for $m>n$ by something that depends on $n$ alone. For the partial sums
$s_n=\sum_{k=1}^n\frac1{k^2}$ the telescope tool does it:
\[ |s_m-s_n|=\sum_{k=n+1}^m\frac1{k^2}\le\sum_{k=n+1}^m\Bigl(\frac1{k-1}-\frac1k\Bigr)=\frac1n-\frac1m<\frac1n . \]
So $N>\frac1\eps$ wins, and $\sum\frac1{k^2}$ converges without anyone knowing its value.

\trap{``It obviously climbs.'' Compute $a_{n+1}-a_n$ and factor it, or induct.}
\trap{Cauchy concerns \emph{all} pairs $m,n\ge N$, not neighbours: $\sqrt{n+1}-\sqrt n\to0$, yet $\sqrt n\to\infty$.}

\begin{yourturn}
\task{t:rec}$a_1=0$, $a_{n+1}=\sqrt{6+a_n}$. Convergent? Limit?
\task{t:heron}$a>0$, $x_1>0$, $x_{n+1}=\frac12\bigl(x_n+\frac a{x_n}\bigr)$: your Übung~3/7, cold.
\task{t:cauchy}Show directly that $s_n=\sum_{k=0}^n2^{-k}$ is a Cauchy sequence.
\end{yourturn}

% =====================================================================================
\section{Series}

\opening{A}{series is not} a new object. It is the sequence of its partial sums, $S_k=\sum_{n=0}^ka_n$, and it converges when
that sequence does. So every series question is a sequence question in disguise, and the tests of chapter~3 are Part~IV's
theorems in working clothes.
\aside{Euler summed $1-1+1-\cdots$ to $\frac12$, and by his rules he was not wrong. Yours are stricter.}

\subsection{The terms of chapter 3}

\entry{Series}{Reihe}
the sequence $(S_k)$ of partial sums; $\sum_{n=0}^\infty a_n$ also names its limit when it exists.

\entry{Geometric series}{geometrische Reihe}
$\sum_{n=0}^\infty x^n=\frac1{1-x}$ for $|x|<1$; divergent for $|x|\ge1$, because then $x^n\not\to0$. Shifted:
$\sum_{n=n_0}^\infty x^n=\frac{x^{n_0}}{1-x}$.

\entry{Harmonic series}{harmonische Reihe}
$\sum\frac1n$ diverges, although $\frac1n\to0$. $\sum\frac1{n^\alpha}$ converges for $\alpha>1$ and diverges for $\alpha\le1$
(compare with $\frac1n$).

\entry{Absolute and conditional convergence}{absolute, bedingte Konvergenz}
absolutely: $\sum|a_n|$ converges, which implies convergence. Conditionally: convergent, not absolutely.
\eg $\sum\frac{(-1)^{n-1}}n$ is conditionally convergent.

\entry{Cauchy product}{Cauchy-Produkt}
$c_n=\sum_{k=0}^na_kb_{n-k}$. If $\sum a_n$ and $\sum b_n$ converge \emph{absolutely}, then $\sum c_n$ converges absolutely to
$\bigl(\sum a_n\bigr)\bigl(\sum b_n\bigr)$.
\eg $\bigl(\sum x^n\bigr)^2=\sum(n+1)x^n=\frac1{(1-x)^2}$ for $|x|<1$; $\exp(x)\exp(y)=\exp(x+y)$.

\entry{Exponential function}{Exponentialfunktion}
$\exp x=\sum_{n=0}^\infty\frac{x^n}{n!}$, absolutely convergent for every $x$; $e=\exp1>\frac52$.

% ------------------------------------------------------------------------------------
\subsection{Move 07\quad Series through their partial sums}
\when{``Bestimmen Sie den Wert der Reihe''; rational terms with a product in the denominator; a geometric series with a shifted
index or a constant factor.}

\begin{panel}[unbreakable]{Skeleton}{telescoping}{Teleskopsumme}
\begin{enumerate}
  \item Partial fractions: $\dfrac1{(n+a)(n+b)}=\dfrac1{b-a}\Bigl(\dfrac1{n+a}-\dfrac1{n+b}\Bigr)$.
  \item Write $S_N$ out, first and last few terms; cross out what cancels.
  \item $S_N=$ (finite leftover) $-$ (terms containing $N$) $\to$ the value.
  \item Geometric: pull out the constant until you see $c\sum_{n=n_0}^\infty q^n=c\,\frac{q^{n_0}}{1-q}$.
\end{enumerate}
\end{panel}

\begin{example}{a telescope with a gap of two}
$\displaystyle\sum_{n=1}^\infty\frac1{n(n+2)}$. Partial fractions: $\frac1{n(n+2)}=\frac12\bigl(\frac1n-\frac1{n+2}\bigr)$.
\[
S_N=\frac12\Bigl[\Bigl(1-\frac13\Bigr)+\Bigl(\frac12-\frac14\Bigr)+\Bigl(\frac13-\frac15\Bigr)+\cdots
+\Bigl(\frac1N-\frac1{N+2}\Bigr)\Bigr].
\]
Every $\frac1k$ with $3\le k\le N$ appears once with $+$ and once with $-$. Two terms survive at each end:
\[ S_N=\frac12\Bigl(1+\frac12-\frac1{N+1}-\frac1{N+2}\Bigr)\ \longrightarrow\ \frac34 . \]
\end{example}

\begin{example}{a geometric series in disguise}
$\displaystyle\sum_{n=2}^\infty\frac{3^n}{4^{n+1}}=\frac14\sum_{n=2}^\infty\Bigl(\frac34\Bigr)^n
=\frac14\cdot\frac{(3/4)^2}{1-3/4}=\frac9{16}$.
\end{example}

\begin{yourturn}
\task{t:tel}$\displaystyle\sum_{n=1}^\infty\frac1{4n^2-1}$.
\task{t:geo}$\displaystyle\sum_{n=0}^\infty\frac{2^n+3^n}{6^n}$.
\task{t:log}Show that $\sum_{n=1}^\infty\log\bigl(1+\frac1n\bigr)$ diverges, although its terms tend to $0$.
\end{yourturn}

% ------------------------------------------------------------------------------------
\subsection{Move 08\quad Choosing a test}
\when{``Untersuchen Sie auf Konvergenz''; with a parameter, ``Für welche $x\in\R$ konvergiert \dots?''}

\begin{panel}{Skeleton}{the order of questions}{Kriterium wählen}
\begin{enumerate}
  \item \textbf{Do the terms tend to $0$?} If not, divergent \de{(notwendige Bedingung)}. Done.
  \item \textbf{Alternating signs?} Test $\sum|a_n|$ first (Move~09). Leibniz only if that fails.
  \item \textbf{Read the shape of $a_n$} and pick from the table.
\end{enumerate}
\medskip
\begin{twocol}{40mm}{Shape of the term}{Reach for, with a one-line example}
\row{$n!$ or $c^n$}{ratio test \de{(Quotientenkriterium)}: $\frac{2^n}{n!}$ has ratio $\frac2{n+1}\le\frac12$ for $n\ge3$}
\row{everything under one $n$-th power}{root test \de{(Wurzelkriterium)}: $\bigl(\frac n{2n+1}\bigr)^n$ has root
$\frac n{2n+1}<\frac12$}
\row{a rational or root expression in $n$}{compare with $\frac1{n^\alpha}$, $\alpha=$ degree below minus degree above:
$\frac n{n^3+1}\le\frac1{n^2}$; \ $\frac1{\sqrt{n^2+1}}\ge\frac1{\sqrt2\,n}$ diverges}
\row{a product in the denominator}{partial sums (Move~07)}
\row{a bounded factor times a summable one}{majorant: $\bigl|\frac{\cos n}{n^2}\bigr|\le\frac1{n^2}$}
\end{twocol}
\end{panel}

\begin{panel}[unbreakable]{Three remarks}{the Skript's exact tools}{}
\begin{itemize}
  \item \textbf{The ratio test needs a fixed $\theta<1$}: $\bigl|\frac{a_{n+1}}{a_n}\bigr|\le\theta$ for all $n\ge n_0$. If you
        computed a limit $\alpha<1$ of the ratios, convert it in one sentence:
        \[ \text{``Da }\Bigl|\frac{a_{n+1}}{a_n}\Bigr|\to\alpha<1,\text{ gilt ab einem }n_0\text{: }
           \Bigl|\frac{a_{n+1}}{a_n}\Bigr|\le\theta:=\frac{1+\alpha}2<1.\text{''} \]
        That is the definition of the limit, played with $\eps=\frac{1-\alpha}2$.
  \item \textbf{Divergence by ratio.} If $\bigl|\frac{a_{n+1}}{a_n}\bigr|\ge1$ from some $n_0$ on, then
        $|a_n|\ge|a_{n_0}|>0$. The terms do not tend to $0$.
  \item \textbf{Minorant.} The Skript names no such criterion, but it is the majorant criterion read backwards. If
        $0\le c_n\le|a_n|$ and $\sum c_n$ diverges, then $\sum|a_n|$ diverges: otherwise it would be a convergent majorant
        of $\sum c_n$.
\end{itemize}
\end{panel}

\begin{example}{factorials over powers, with no $e$ in sight}
$\displaystyle\sum_{n=1}^\infty\frac{n!}{n^n}$. The ratio:
\[ \frac{a_{n+1}}{a_n}=\frac{(n+1)!}{(n+1)^{n+1}}\cdot\frac{n^n}{n!}=\Bigl(\frac n{n+1}\Bigr)^n=\frac1{(1+\frac1n)^n}. \]
Bernoulli: $(1+\frac1n)^n\ge1+n\cdot\frac1n=2$. So the ratio is at most $\theta=\frac12$ for every $n$: absolutely convergent.
\end{example}

\begin{example}{a parameter, and the edges by hand}
For which $x\in\R$ does $\sum_{n=1}^\infty\frac{x^n}n$ converge? For $x\ne0$ the ratio of absolute values is
$|x|\frac n{n+1}$.
\begin{itemize}
  \item $|x|<1$: the ratio is $\le\theta=|x|<1$. Absolutely convergent.
  \item $|x|>1$: the ratio is $\ge1$ once $n\ge\frac1{|x|-1}$. The terms do not tend to $0$: divergent.
  \item $x=1$: harmonic, divergent. \ $x=-1$: Leibniz, convergent (conditionally).
\end{itemize}
Answer: exactly for $x\in[-1,1)$.
\end{example}

\trap{$|a_{n+1}/a_n|<1$ for every $n$ is not enough. For $\frac1n$ the ratio is $\frac n{n+1}<1$, and the series diverges. You
need one fixed $\theta<1$.}
\trap{A ratio or root limit equal to $1$ decides nothing: $\sum\frac1n$ and $\sum\frac1{n^2}$ both have limit $1$.}
\trap{The edges of a parameter range are separate questions. Test them by hand.}

\begin{yourturn}
\task{t:ratio}$\displaystyle\sum\frac{n^2}{3^n}$.
\task{t:root}$\displaystyle\sum\Bigl(\frac n{3n-1}\Bigr)^n$.
\task{t:comp2}$\displaystyle\sum\frac1{\sqrt{n(n+1)}}$.
\task{t:hard}$\displaystyle\sum\frac{3^n\,n!}{n^n}$. Use $(1+\frac1n)^n\le3$.
\task{t:param}For which $x\in\R$ does $\sum_{n=1}^\infty n\,x^n$ converge?
\end{yourturn}

% ------------------------------------------------------------------------------------
\subsection{Move 09\quad Absolute or conditional}
\when{``Untersuchen Sie auf Konvergenz und absolute Konvergenz''; any alternating series.}

\begin{panel}[unbreakable]{Skeleton}{absolute or conditional}{absolut oder bedingt}
\begin{enumerate}
  \item \textbf{Absolute first.} Test $\sum|a_n|$ with Move~08. If it converges: absolutely convergent, hence convergent. Done.
  \item \textbf{Leibniz.} If $\sum|a_n|$ diverges, write $a_n=(-1)^nb_n$ and check all three conditions: $b_n\ge0$,
        $b_n\to0$, $(b_n)$ decreasing. Prove the last one, by cross-multiplying; no derivatives.
  \item \textbf{Verdict:} convergent, not absolutely: \de{bedingt konvergent}.
\end{enumerate}
\end{panel}

\begin{example}{conditional convergence, monotonicity by algebra}
$\displaystyle\sum_{n=1}^\infty(-1)^n\frac n{n^2+1}$.

\emph{Absolute?} $\frac n{n^2+1}\ge\frac n{2n^2}=\frac1{2n}$. A harmonic minorant: $\sum|a_n|$ diverges.

\emph{Leibniz.} $b_n=\frac n{n^2+1}\ge0$ and $b_n\le\frac1n\to0$. Decreasing:
\[\begin{aligned}
b_{n+1}\le b_n&\iff(n+1)(n^2+1)\le n\bigl((n+1)^2+1\bigr)\\
&\iff n^3+n^2+n+1\le n^3+2n^2+2n\\
&\iff 1\le n^2+n,
\end{aligned}\]
true for $n\ge1$. So the series converges, but not absolutely.
\end{example}

\why{Why does the Cauchy product need absolute convergence?}{Regrouping the products $a_kb_\ell$ along diagonals is harmless
only when their absolute values are summable. Without it, square $\sum\frac{(-1)^n}{\sqrt{n+1}}$, which converges by Leibniz.
Since $(k+1)(n-k+1)\le\bigl(\frac{n+2}2\bigr)^2$,
\[ |c_n|=\sum_{k=0}^n\frac1{\sqrt{(k+1)(n-k+1)}}\ge\sum_{k=0}^n\frac2{n+2}=\frac{2(n+1)}{n+2}\ge1 . \]
The terms do not tend to $0$: the product of two convergent series diverges.}

\begin{yourturn}
\task{t:leib1}$\displaystyle\sum\frac{(-1)^n}{2n+1}$.
\task{t:leib2}$\displaystyle\sum(-1)^n\frac{n^2}{2^n}$.
\task{t:leib3}$\displaystyle\sum(-1)^n\bigl(\sqrt{n+1}-\sqrt n\bigr)$.
\end{yourturn}

% =====================================================================================
\section{Functions and continuity}

\opening{C}{hapter 4 asks} one question in many costumes: if $x$ stays close to $x_0$, does $f(x)$ stay close to $f(x_0)$? The
Skript answers it twice, with sequences and with $\eps$ and $\delta$, and proves the two answers agree. Use whichever is shorter.
Sequences are usually shorter for \emph{dis}continuity; $\eps$-$\delta$ is what the question asks for when it says so.

\subsection{The terms of chapter 4}

\entry{Accumulation point of a set}{Häufungspunkt einer Menge}
$y$ is one for $D$ if some sequence in $D\setminus\{y\}$ tends to $y$. A contact point \de{(Berührpunkt)} only needs a sequence
in $D$. \eg $0$ for $D=(0,1)$, although $0\notin D$.

\entry{Limit of a function}{Grenzwert einer Funktion}
$\lim_{x\to y}f(x)=c$ iff $f(x_n)\to c$ for \emph{every} sequence $x_n\to y$ in $D\setminus\{y\}$. Equivalently (a theorem):
$\forall\eps>0\ \exists\delta>0\ \forall x\in D: 0<|x-y|<\delta\Rightarrow|f(x)-c|<\eps$. One-sided: only $x_n>y$ (written
$x\downarrow y$) or $x_n<y$ ($x\uparrow y$). At $\infty$: every sequence $x_n\to\infty$.

\entry{Continuity}{Stetigkeit}
$f$ is continuous at $y\in D$ iff $\forall\eps\ \exists\delta\ \forall x\in D: |x-y|<\delta\Rightarrow|f(x)-f(y)|<\eps$;
equivalently $x_n\to y\Rightarrow f(x_n)\to f(y)$. $C(D)$ is the set of functions continuous on $D$.
\eg polynomials, $\exp$, $\sin$, $\cos$, $|x|$. \nonex the Dirichlet function $D$, $1$ on $\Q$ and $0$ elsewhere: continuous
nowhere.

\entry{Uniform continuity}{gleichmäßige Stetigkeit}
$\forall\eps\ \exists\delta\ \forall x,y\in D: |x-y|<\delta\Rightarrow|f(x)-f(y)|<\eps$. The $\delta$ may depend on $\eps$
only. \eg $x$ on $\R$; every continuous $f$ on $[a,b]$. \nonex $x^2$ on $\R$; $\frac1x$ on $(0,1]$.

\entry{Continuous extension}{stetige Fortsetzung}
a continuous $\tilde f$ on a larger set that agrees with $f$ on $D$. At a single missing point $y$ it exists iff
$\lim_{x\to y}f(x)$ exists; then $\tilde f(y)$ is that limit. \eg $\frac{x^2-4}{x-2}$ extends by the value $4$ at $2$.
\nonex $\frac1x$ at $0$.

\entry{Step function}{Treppenfunktion}
$\varphi\colon[a,b]\to\R$ with a partition $a=t_0<\dots<t_n=b$ such that $\varphi$ is constant on each $(t_{k-1},t_k)$; the
values at the $t_k$ are free. \eg $0$ on $[0,1)$, $1$ on $[1,2]$. \nonex $x$ on $[0,1]$.

\entry{Image of a compact interval}{Bild kompakter Intervalle}
for continuous $f$: $f([a,b])=[\min f,\max f]$, a closed bounded interval. The extreme value theorem gives the ends, the
intermediate value theorem everything between.

\entry{Power, root, logarithm}{Potenz, Wurzel, Logarithmus}
$\sqrt[m]a$ is the unique $x>0$ with $x^m=a$; $a^{p/q}=\sqrt[q]{a^p}$. $\log=\exp^{-1}\colon(0,\infty)\to\R$;
$a^x=\exp(x\log a)$ for $a>0$; $\log_a$ is the inverse of $x\mapsto a^x$ ($a\ne1$), and $\log_ay=\frac{\log y}{\log a}$.

\entry{Sine and cosine}{Sinus, Kosinus}
defined by their series; addition theorems; $\sin(-x)=-\sin x$, $\cos(-x)=\cos x$; $\lim_{x\to0}\frac{\sin x}x=1$. One line
gives $\sin^2x+\cos^2x=1$, hence $|\sin x|\le1$:
\[ 1=\cos0=\cos(x-x)=\cos x\cos(-x)-\sin x\sin(-x)=\cos^2x+\sin^2x . \]

% ------------------------------------------------------------------------------------
\subsection{Move 10\quad Exponential and logarithm}
\when{equations or inequalities with $\exp$, $\log$, $a^x$; limits involving $\exp$.}

\begin{panel}[unbreakable]{Rules}{everything you may use}{Rechenregeln}
\begin{itemize}
  \item $\exp(x+y)=\exp x\exp y$, \ $\exp(-x)=\frac1{\exp x}$, \ $\exp x>0$; strictly increasing, bijective $\R\to(0,\infty)$.
  \item $\log(xy)=\log x+\log y$, \ $\log1=0$, \ $\log e=1$; strictly increasing and continuous.
  \item $a^x=\exp(x\log a)$, so $a^{x+y}=a^xa^y$ and $\log(a^x)=x\log a$.
  \item $\exp x\ge1+x$ for $x\ge0$; \ $e>\frac52$; \ $|\exp x-1|\le2|x|$ for $|x|\le1$.
  \item Several of $\exp(x)$, $\exp(2x)$, $\exp(-x)$? Substitute $t=\exp x>0$ (Vorbereitung~§10).
\end{itemize}
\end{panel}

\begin{example}{three small ones}
\emph{Solve $2^{x+1}=3^x$.} $\log$ is strictly increasing, so taking it is an equivalence:
\[\begin{aligned}
(x+1)\log2=x\log3&\iff x(\log3-\log2)=\log2\\
&\iff x=\frac{\log2}{\log\frac32}.
\end{aligned}\]
\emph{Solve $\log x+\log(x-1)=\log6$.} Domain first: $x>1$. There, $\log\bigl(x(x-1)\bigr)=\log6\iff x^2-x-6=0\iff
x\in\{3,-2\}$. Only $x=3$ lies in the domain.

\emph{Show $\lim_{x\to\infty}\frac x{\exp x}=0$.} For $x>0$, $\exp x\ge1+x+\frac{x^2}2>\frac{x^2}2$, so
$0<\frac x{\exp x}<\frac2x\to0$.
\end{example}

\begin{yourturn}
\task{t:exp1}Solve $\exp x+\exp(-x)=4$.
\task{t:exp2}$\displaystyle\lim_{x\to\infty}\frac{\exp x-\exp(-x)}{\exp x+\exp(-x)}$ (your Übung~5/2b, cold).
\task{t:exp3}$\displaystyle\lim_{x\to\infty}\frac{x^2}{\exp x}$, with Skript tools only.
\end{yourturn}

% ------------------------------------------------------------------------------------
\subsection{\texorpdfstring{Move 11\quad Continuity and limits with $\eps$ and $\delta$}{Move 11: continuity and limits with epsilon and delta}}
\when{``Zeigen Sie mit dem $\eps$-$\delta$-Kriterium, dass $f$ in $x_0$ stetig ist'' or ``dass $\lim_{x\to x_0}f(x)=c$''.}

\begin{panel}[unbreakable]{Skeleton}{epsilon-delta}{Epsilon-Delta}
\begin{twodocs}
\begin{enumerate}
  \item Factor $|x-x_0|$ out: $|f(x)-f(x_0)|=|x-x_0|\cdot|h(x)|$.
  \item Restrict $\delta\le1$ (or $\delta\le\frac{x_0}2$ if $f$ blows up at $0$). Inside the window, bound $|h(x)|\le C$.
  \item Then $C\delta\le\eps$ is enough: $\delta\le\frac\eps C$.
\end{enumerate}
&
``Sei $\eps>0$. Setze $\delta:=\min\{1,\frac\eps C\}$. Für $x\in D$ mit $|x-x_0|<\delta$ ist $|x-x_0|<1$, also $|h(x)|\le
C$, und
\[ |f(x)-f(x_0)|\le C|x-x_0|<C\delta\le\eps .\text{''} \]
\end{twodocs}
For a limit, the window is $0<|x-x_0|<\delta$: simplify first, using $x\ne x_0$.
\end{panel}

\begin{example}{the square at an arbitrary point}
$f(x)=x^2$ is continuous at every $x_0\in\R$. Factor: $|x^2-x_0^2|=|x-x_0|\,|x+x_0|$. If $|x-x_0|<1$, then
\[ |x+x_0|\le|x-x_0|+2|x_0|<1+2|x_0|=:C . \]
So $\delta:=\min\bigl\{1,\frac\eps{1+2|x_0|}\bigr\}$ works: $|x-x_0|<\delta\Rightarrow|x^2-x_0^2|<C\delta\le\eps$.
\end{example}

\begin{example}{the reciprocal, where the danger is at $0$}
$f(x)=\frac1x$ on $(0,\infty)$, at $x_0>0$. Here $\bigl|\frac1x-\frac1{x_0}\bigr|=\frac{|x-x_0|}{x\,x_0}$, which is large
when $x$ is near $0$. Restrict $\delta\le\frac{x_0}2$: then $x>\frac{x_0}2$ and
\[ \Bigl|\frac1x-\frac1{x_0}\Bigr|<\frac{2}{x_0^2}\,|x-x_0| . \]
So $\delta:=\min\bigl\{\frac{x_0}2,\frac{x_0^2\eps}2\bigr\}$. This was your Übung~5/3a. Look at how $\delta$ depends on $x_0$.
\end{example}

\par\needspace{95mm}
\begin{center}
\begin{tikzpicture}
\begin{axis}[klplot,width=.88\linewidth,height=62mm,xmin=0,xmax=2.2,ymin=0,ymax=5.4,
  xlabel={$x$},ylabel={$\frac1x$},xtick={0.25,1,2},xticklabels={$\frac14$,$1$,$2$},ytick={1,4},ylabel style={rotate=-90}]
  \fill[pattern={dither25},pattern color=klink] (axis cs:0,0.5) rectangle (axis cs:2.2,1.5);
  \fill[pattern={dither25},pattern color=klink] (axis cs:0,3.5) rectangle (axis cs:2.2,4.5);
  \fill[pattern={dither50},pattern color=klink] (axis cs:0.6667,0) rectangle (axis cs:1.3333,0.3);
  \fill[pattern={dither50},pattern color=klink] (axis cs:0.2222,0) rectangle (axis cs:0.2778,0.3);
  \addplot[klhi,domain=0.185:2.2,samples=200] {1/x};
  \draw[kldim,dashed] (axis cs:1,0)--(axis cs:1,1);
  \draw[kldim,dashed] (axis cs:0.25,0)--(axis cs:0.25,4);
  \node[anchor=west,text=klink] at (axis cs:1.36,0.17) {$\delta=\frac13$};
  \node[anchor=west,text=klink] at (axis cs:0.30,0.17) {$\delta=\frac1{36}$};
\end{axis}
\end{tikzpicture}
\end{center}
\plotcaption{The same $\eps=\frac12$ at two points: the bands are $f(x_0)\pm\eps$, the strips on the axis are the windows
$x_0\pm\delta$ that land inside them. At $x_0=1$ any $\delta\le\frac13$ works; at $x_0=\frac14$ you need $\delta\le\frac1{36}$. As
$x_0\to0$ the admissible $\delta$ shrinks to $0$. Continuous at every point, yet no single $\delta$ serves them all: that is the
failure of uniform continuity (Move~15).}

\noindent\textbf{Discontinuity and non-existence, with sequences.} One sequence $x_n\to x_0$ with $f(x_n)\not\to f(x_0)$ kills
continuity; two sequences with different limits of $f$ kill a limit. The Dirichlet function at any $x_0$: rationals
$q_n\to x_0$ and irrationals $r_n\to x_0$ exist by density, and $D(q_n)=1$, $D(r_n)=0$.

\trap{$\delta$ depending on $x$, as in $\delta=\frac\eps{|x+x_0|}$. It may depend on $\eps$ and $x_0$, never on $x$.}
\trap{For a limit, forgetting $0<|x-x_0|$. The value at $x_0$ is irrelevant to the limit, and may not exist.}

\begin{yourturn}
\task{t:ed1}$f(x)=x^3$ is continuous at $x_0=1$, with $\eps$ and $\delta$.
\task{t:ed2}$\displaystyle\lim_{x\to2}\frac{x^2-4}{x-2}=4$, with $\eps$ and $\delta$.
\task{t:ed3}$f(x)=\frac1x$ is continuous at $x_0=2$, with $\eps$ and $\delta$.
\end{yourturn}

% ------------------------------------------------------------------------------------
\subsection{Move 12\quad One-sided limits, limits at infinity, glue points}
\when{piecewise definitions; $|x|$ or the Gauß bracket $[x]$; ``Für welches $a$ ist $f$ stetig?''; $x\to\pm\infty$.}

\begin{panel}[unbreakable]{Skeleton}{one-sided limits}{einseitige Grenzwerte}
\begin{enumerate}
  \item From one side you know the sign of $x-x_0$. Use it to resolve $|x|$, $[x]$, or the case split.
  \item The two-sided limit exists iff both one-sided limits exist and agree. At a point of $D$, continuity means both equal
        $f(x_0)$.
  \item At $\infty$: $\forall\eps\ \exists M: x>M\Rightarrow|f(x)-c|<\eps$. Estimate as in Move~04, with $x$ for $n$.
\end{enumerate}
\end{panel}

\begin{example}{choose the parameter}
For which $a\in\R$ is $f$ continuous on $\R$?
\[ f(x)=\begin{cases}\dfrac{x^2-1}{x-1}, & x<1,\\[6pt] ax+1, & x\ge1.\end{cases} \]
Away from $1$, each piece is continuous by the rules (the denominator $x-1$ does not vanish for $x<1$). At $1$: from the left,
$f(x)=x+1\to2$. From the right and at the point, $ax+1\to a+1=f(1)$. Continuous at $1$ iff $2=a+1$. So $a=1$.
\end{example}

\begin{yourturn}
\task{t:ext}$f(x)=\dfrac{\sqrt{1+x}-1}{x}$ on $[-1,0)\cup(0,\infty)$. Which value at $0$ makes $f$ continuous?
\task{t:inf}$\displaystyle\lim_{x\to\infty}\frac{2x^2+x}{x^2-3}=2$, from the definition (find $M$ for each $\eps$).
\par\smallskip\noindent Then your Übung~5/4 ($[x]$ and $[x]+\sqrt{x-[x]}$ on $(0,2)$) cold.
\end{yourturn}

% ------------------------------------------------------------------------------------
\subsection{Move 13\quad Continuity by the rules}
\when{``Zeigen Sie, dass $f$ stetig ist'' for a formula assembled from known functions.}

\begin{panel}[unbreakable]{Skeleton}{the rules}{Rechenregeln}
\begin{enumerate}
  \item Name the bricks: polynomials, $\exp$, $\log$, $\sin$, $\cos$, $|\cdot|$, $\sqrt{\cdot}$ (continuous on $[0,\infty)$ as an
        inverse function).
  \item Name the rule at each joint: sum, product, quotient (check that the denominator never vanishes), composition.
  \item Where no rule applies (glue points, zeros of a denominator): one-sided limits, or a squeeze
        $|f(x)-f(x_0)|\le g(x)\to0$.
\end{enumerate}
\end{panel}

\begin{example}{a squeeze at the glue point}
$f(x)=x\sin\frac1x$ for $x\ne0$, $f(0)=0$, is continuous on $\R$. On $\R\setminus\{0\}$: $\frac1x$ is a quotient with nonzero
denominator, $\sin$ is continuous, so is the composition, and so is the product with $x$. At $0$: $|\sin t|\le1$ for all $t$,
so
\[ |f(x)-f(0)|=|x|\,\Bigl|\sin\frac1x\Bigr|\le|x|<\eps\quad\text{for }|x|<\delta:=\eps . \]
\end{example}

\trap{``Quotient of continuous functions'' without checking that the denominator never vanishes on the domain.}
\trap{``Composition of continuous functions'' for a piecewise $f$. The pieces are covered; the glue point is not.}

\begin{yourturn}
\task{t:sq}$f(x)=\sqrt x\,\sin\frac1x$ on $(0,\infty)$. Can it be continuously extended to $[0,\infty)$? (Übung~5/8b, cold.)
\task{t:dir}$g(x)=x\,D(x)$, with $D$ the Dirichlet function, is continuous at $0$ and nowhere else.
\end{yourturn}

% ------------------------------------------------------------------------------------
\subsection{Move 14\quad The intermediate value theorem}
\when{``Zeigen Sie, dass die Gleichung \dots\ (mindestens / genau) eine Lösung in \dots\ besitzt''; fixed points; polynomials of
odd degree.}

\begin{panel}[unbreakable]{Skeleton}{zeros}{Nullstellen}
\begin{enumerate}
  \item Move everything to one side: $h=f-g$, or $h=f-c$, or $h(x)=f(x)-x$ for a fixed point.
  \item $h$ is continuous on the \emph{closed} interval $[a,b]$: say why.
  \item Compute $h(a)$ and $h(b)$: opposite signs, or one of them $0$.
  \item IVT: some $p\in[a,b]$ has $h(p)=0$; with strict signs, $p\in(a,b)$.
  \item ``Genau eine'': $h$ strictly monotone, shown by algebra. ``Mindestens drei'': three sign changes on disjoint intervals.
\end{enumerate}
\end{panel}

\begin{example}{existence and uniqueness}
$\exp x=3-x$ has exactly one real solution, and it lies in $(0,1)$. Let $h(x)=\exp x+x-3$, continuous on $\R$ as a sum.
\[ h(0)=1-3=-2<0,\qquad h(1)=e-2>\tfrac52-2>0 . \]
The IVT gives a zero in $(0,1)$. Unique: $\exp$ and $x\mapsto x$ are strictly increasing, so $h$ is strictly increasing on $\R$
and takes the value $0$ at most once.
\end{example}

\par\needspace{85mm}
\begin{center}
\begin{tikzpicture}
\begin{axis}[klplot,width=.88\linewidth,height=60mm,xmin=0,xmax=1.12,ymin=-2.05,ymax=1.25,
  xlabel={$x$},ylabel={$f$},xtick={0,0.5,0.625,0.75,1},xticklabels={$0$,$\frac12$,$\frac58$,$\frac34$,$1$},ytick={-1,0,1},
  ylabel style={rotate=-90}]
  \draw[kldim,dashed] (axis cs:0,0)--(axis cs:1.05,0);
  \addplot[klhi,domain=0:1,samples=100] {x^3+x-1};
  \draw[klink,line width=1.3pt] (axis cs:0,-1.3)--(axis cs:1,-1.3) node[right]{$I_0$};
  \draw[klink,line width=1.3pt] (axis cs:0.5,-1.5)--(axis cs:1,-1.5) node[right]{$I_1$};
  \draw[klink,line width=1.3pt] (axis cs:0.5,-1.7)--(axis cs:0.75,-1.7) node[right]{$I_2$};
  \draw[klink,line width=1.3pt] (axis cs:0.625,-1.9)--(axis cs:0.75,-1.9) node[right]{$I_3$};
  \fill[klhi] (axis cs:0.6823,0) circle (1.7pt);
  \node[anchor=south east,text=klhi] at (axis cs:0.6823,0.02) {$p$};
\end{axis}
\end{tikzpicture}
\end{center}
\plotcaption{The proof of the IVT is a hunt by halves, here for $f(x)=x^3+x-1$ on $[0,1]$. Keep the half where $f$ is $\le0$ on
the left end and $\ge0$ on the right. The left ends climb, the right ends fall, both converge (monotone convergence) to one
$p$, and continuity squeezes $f(p)\le0\le f(p)$.}

\trap{Applying the IVT on an open interval, or on a set with a hole. It needs $[a,b]$ and continuity on all of it.
$x^2-2$ changes sign on $[0,2]\cap\Q$ and has no zero there.}
\trap{Existence is all the IVT gives. Uniqueness is a separate argument.}

\begin{yourturn}
\task{t:ivt3}$x^5-3x+1=0$ has at least three real solutions.
\task{t:chord}$f\colon[0,1]\to\R$ continuous with $f(0)=f(1)$. Show: some $x\in[0,\frac12]$ has $f(x)=f(x+\frac12)$.
\task{t:odd}Every polynomial of odd degree has a real zero (your Übung~5/5a, cold).
\end{yourturn}

% ------------------------------------------------------------------------------------
\subsection{Move 15\quad Uniform continuity}
\when{``Ist $f$ gleichmäßig stetig auf $D$?'' On a closed bounded interval the answer is always yes (Heine--Cantor). The
question is only interesting on open or unbounded sets.}

\begin{panel}[unbreakable]{Skeleton}{uniform continuity}{gleichmäßige Stetigkeit}
\textbf{To show it:}
\begin{enumerate}
  \item $D=[a,b]$ and $f$ continuous: Heine--Cantor. Done.
  \item Lipschitz: $|f(x)-f(y)|\le L|x-y|$ for all $x,y\in D$. Then $\delta:=\frac\eps L$.
  \item Otherwise find $|f(x)-f(y)|\le\varphi(|x-y|)$ with $\varphi(t)\to0$. Example: $|\sqrt x-\sqrt y|\le\sqrt{|x-y|}$, so
        $\delta:=\eps^2$.
\end{enumerate}
\textbf{To refute it:}
\begin{enumerate}
  \setcounter{enumi}{3}
  \item Two sequences in $D$ with $|x_n-y_n|\to0$ but $|f(x_n)-f(y_n)|\ge\eps_0$. Look where $f$ is steepest: at $\infty$ for
        $x^2$, near $0$ for $\frac1x$.
  \item Shortcut: uniformly continuous on a bounded $(a,b)$ implies bounded, because $f$ then extends continuously to $[a,b]$.
        So a continuous, unbounded $f$ on a bounded open interval is not uniformly continuous.
\end{enumerate}
\end{panel}

\begin{example}{the square on a half-line, refuted}
$x^2$ on $[0,\infty)$: take $x_n=n$, $y_n=n+\frac1n$. Then $|x_n-y_n|=\frac1n\to0$, but
\[ y_n^2-x_n^2=2+\frac1{n^2}\ge2 . \]
With $\eps_0=1$: whatever $\delta>0$ the opponent offers, choose $n>\frac1\delta$. The pair is closer than $\delta$ and its
values differ by more than $1$.
\end{example}

\begin{example}{the root on a half-line, proved}
$\sqrt x$ on $[0,\infty)$ is uniformly continuous. Let $0\le y\le x$. Since $\sqrt x\sqrt y\ge\sqrt y\sqrt y=y$,
\[ \bigl(\sqrt x-\sqrt y\bigr)^2=x-2\sqrt x\sqrt y+y\le x-2y+y=x-y . \]
So $|\sqrt x-\sqrt y|\le\sqrt{|x-y|}$, and $\delta:=\eps^2$ serves every point. Yet $\sqrt x$ is not Lipschitz:
$\frac{|\sqrt x-\sqrt0|}{|x-0|}=\frac1{\sqrt x}$ is unbounded. Steep at one point is allowed, as long as it is steep in a bounded
way.
\end{example}

\why{Why does Heine--Cantor need a closed bounded interval?}{The proof assumes bad pairs $x_n,y_n$ and needs a subsequence that
converges to a point \emph{of the domain}, where continuity can be used. Bolzano--Weierstraß supplies the subsequence only for
bounded sequences, and the limit stays in the domain only if the interval is closed. Drop ``bounded'': $x^2$ on $[0,\infty)$.
Drop ``closed'': $\frac1x$ on $(0,1]$.}

\begin{yourturn}
\task{t:uc1}$\frac1{1+x^2}$ is uniformly continuous on $\R$: show it is Lipschitz with $L=1$.
\task{t:uc2}$\frac1x$ is not uniformly continuous on $(0,1]$ (Übung~5/3b, cold), but it is on $[1,\infty)$.
\task{t:uc3}$\exp$ is not uniformly continuous on $\R$.
\end{yourturn}

% ------------------------------------------------------------------------------------
\subsection{Move 16\quad Inverting strictly monotone functions}
\when{``Zeigen Sie, dass $f$ bijektiv ist und $f^{-1}$ stetig''; ``Bestimmen Sie $f^{-1}(c)$''; roots, $\log$, $a^x$.}

\begin{panel}[unbreakable]{Skeleton}{the inverse chain}{Umkehrfunktion}
\begin{enumerate}
  \item \textbf{Continuous:} by the rules.
  \item \textbf{Strictly monotone:} $f(y)-f(x)>0$ for $x<y$, factored, or as a sum of increasing pieces. Hence injective.
  \item \textbf{Image:} values or limits at the ends of the interval, then the IVT: $f(I)$ is the interval between them.
  \item \textbf{The theorem} \de{(Satz über die Umkehrfunktion)}: $f\colon I\to f(I)$ is bijective, and $f^{-1}$ is continuous
        and strictly monotone in the same direction.
  \item \textbf{Values of $f^{-1}$:} guess $x$ with $f(x)=c$. Injectivity makes the guess the answer.
\end{enumerate}
\end{panel}

\begin{example}{the chain, run once}
$f\colon\R\to\R$, $f(x)=x+\exp x$.
\begin{enumerate}
  \item Continuous, as a sum.
  \item Strictly increasing, as a sum of strictly increasing functions.
  \item Image. Let $c\in\R$, $a=-|c|-1<0$, $b=|c|+1>0$. Since $0<\exp a<1$, $f(a)<a+1=-|c|\le c$. And $f(b)>b>c$. The IVT
        on $[a,b]$ gives $f(x)=c$. So $f(\R)=\R$.
  \item So $f^{-1}\colon\R\to\R$ exists, is continuous and strictly increasing.
  \item $f(0)=0+1=1$, so $f^{-1}(1)=0$.
\end{enumerate}
\end{example}

\begin{yourturn}
\task{t:inv}$f\colon[0,\infty)\to[0,\infty)$, $f(x)=\frac{x^2}{1+x}$. Bijective? Find $f^{-1}$.
\par\smallskip\noindent Then Vorbereitung~§16 cold.
\end{yourturn}

% =====================================================================================
\section{The theorems}

\opening{F}{or every Satz} the test wants two things: the statement with \emph{every} hypothesis, and the proof or its idea.
Below, each theorem in that shape, in the Skript's order. The drill: cover the Idea lines, write them from memory on a blank
page, then compare. A missing hypothesis costs as much as a missing step.

\subsection{Chapter 1}

\satz{Inverses are unique; $0\cdot x=0$}{Eindeutigkeit der Inversen}
In a field the additive and multiplicative inverses are unique, and $0\cdot x=0$, $(-1)x=-x$.
\idea If $y,z$ are both additive inverses of $x$: $y=y+0=y+(x+z)=(y+x)+z=0+z=z$. Then $0x=(0+0)x=0x+0x$; add $-(0x)$. And
$x+(-1)x=(1-1)x=0$, so $(-1)x$ is the inverse of $x$.

\satz{$\sqrt2$ is irrational}{Irrationalität}
No $r\in\Q$ has $r^2=2$.
\idea Suppose $r=\frac pq>0$ in lowest terms. $p^2=2q^2$, so $p^2$ is even, so $p$ is even (an odd $m=2k+1$ has odd square).
Then $p=2k$, $q^2=2k^2$, so $q$ is even too. Contradiction.

\satz{Bernoulli's inequality}{Bernoulli-Ungleichung}
For $h\ge-1$ and $n\in\N$: $(1+h)^n\ge1+nh$.
\idea Induction. The step multiplies the hypothesis by $1+h$, which is allowed because $1+h\ge0$: this is where $h\ge-1$ is
used. Then $(1+nh)(1+h)=1+(n+1)h+nh^2\ge1+(n+1)h$.

\satz{Geometric sum}{geometrische Summenformel}
For $q\ne1$: $\sum_{k=0}^nq^k=\frac{1-q^{n+1}}{1-q}$.
\idea Telescope: $(1-q)\sum_{k=0}^nq^k=\sum_{k=0}^nq^k-\sum_{k=1}^{n+1}q^k=1-q^{n+1}$. Or induction.

\satz{Binomial theorem}{binomischer Lehrsatz}
$(a+b)^n=\sum_{k=0}^n\binom nka^kb^{n-k}$.
\idea Induction. Multiply by $a+b$, shift the index in one sum, and merge with Pascal: $\binom n{k-1}+\binom nk=\binom{n+1}k$.

\satz{Archimedean property}{Archimedische Eigenschaft}
For every $r\in\R$ there is $n\in\N$ with $n>r$.
\idea Otherwise $\N$ is bounded above by $r$, so $s=\sup\N$ exists. $s-1$ is not an upper bound, so some $n>s-1$. Then
$n+1>s$ and $n+1\in\N$. Contradiction.

\satz{$\Q$ is countable}{Abzählbarkeit von $\Q$}
There is a bijection $\N\to\Q$.
\idea List the positive fractions $\frac nm$ along the diagonals $n+m=2,3,4,\dots$, skipping those not in lowest terms. That is a
bijection $\varphi\colon\N\to\Q^+$. Then $0,\ \varphi(1),\ -\varphi(1),\ \varphi(2),\ -\varphi(2),\dots$

\satz{$\R$ is uncountable}{Überabzählbarkeit}
No sequence $x_1,x_2,\dots$ contains every point of $[0,1]$. (On the Erwartungshorizont; not in your Skript excerpt.)
\idea Split $[0,1]$ into three closed thirds. A point lies in at most two of them, so one third $I_1$ misses $x_1$. Split $I_1$
into thirds and pick one, $I_2$, that misses $x_2$; and so on. The nested intervals share a point $x$. It is no $x_n$, because
$x\in I_n$ and $x_n\notin I_n$.

\satz{Existence of the square root}{Existenz der Quadratwurzel}
For $a\ge0$ there is exactly one $x\ge0$ with $x^2=a$.
\idea Unique: $x^2=y^2$ with $x,y\ge0$ gives $(x-y)(x+y)=0$. Exists: $M=\{z\ge0: z^2<a\}$ is nonempty and bounded by $1+a$; let
$s=\sup M$. If $s^2<a$, then $s+\delta\in M$ for small $\delta$, with
$\delta<\min\{1,\frac{a-s^2}{2s+1}\}$: $s$ is no upper bound. If $s^2>a$, then $s-\delta$ is still an upper bound, with
$\delta<\min\{s,\frac{s^2-a}{2s+1}\}$: $s$ is not the least. So $s^2=a$. The $m$-th root goes the same way.

\satz{Nested intervals}{Intervallschachtelungsprinzip}
Closed bounded $I_n=[a_n,b_n]$ with $I_{n+1}\subseteq I_n$ have a common point; exactly one if $b_n-a_n\to0$.
\idea $(a_n)$ increases and every $b_m$ bounds it. $x=\sup_na_n$ satisfies $a_n\le x\le b_n$ for all $n$. Two common points
$x,y$ have $|x-y|\le b_n-a_n\to0$.

\satz{$\Q$ is dense}{Dichtheit von $\Q$}
Between any two reals $x<y$ lie a rational and an irrational number.
\idea Archimedes: $n(y-x)>1$. Then some integer $m$ satisfies $nx<m<ny$, so $x<\frac mn<y$. For an irrational, find a rational
between $x$ and $y-\frac{\sqrt2}M$ and add $\frac{\sqrt2}M$.

\subsection{Chapter 2}

\satz{The limit is unique}{Eindeutigkeit des Grenzwerts}
\idea $a_n\to a$ and $a_n\to b$: $a-b=(a-a_n)+(a_n-b)$ is a constant null sequence, so $a=b$.

\satz{Convergent implies bounded}{konvergent $\Rightarrow$ beschränkt}
\idea Play $\eps=1$: $|a_n|<|a|+1$ for $n\ge N$. Take the maximum with the finitely many $|a_1|,\dots,|a_{N-1}|$.

\satz{Limit laws}{Rechenregeln}
$a_n\to a$, $b_n\to b$ $\Rightarrow$ $a_n+b_n\to a+b$, $ca_n\to ca$, $a_nb_n\to ab$; and $\frac{a_n}{b_n}\to\frac ab$ if
$b\ne0$.
\idea Sums: $\frac\eps2+\frac\eps2$. Products: $a_nb_n-ab=(a_n-a)b_n+a(b_n-b)$ with $(b_n)$ bounded. Quotients: $|b_n|>\frac{|b|}2$
eventually.

\satz{Comparison for null sequences}{Vergleichskriterium}
$|a_n|\le b_n$ for $n\ge N$ and $b_n\to0$ $\Rightarrow$ $a_n\to0$.
\idea Take the larger of the two thresholds: $|a_n|\le b_n<\eps$.

\satz{Monotone convergence}{Monotoniekriterium}
Increasing and bounded above $\Rightarrow$ convergent, with limit $\sup\{a_n\}$.
\idea $a=\sup\{a_n\}$. $a-\eps$ is not an upper bound, so $a_N>a-\eps$; monotone, so $a-\eps<a_n\le a$ for all $n\ge N$.

\satz{Bolzano--Weierstraß}{Satz von Bolzano--Weierstraß}
Every bounded sequence in $\R$ has a convergent subsequence.
\idea Bisect $[A,B]$ and keep a half that contains infinitely many terms. The nested intervals have lengths $2^{-k}(B-A)\to0$ and
shrink to one point $x$. Pick $n_1<n_2<\cdots$ with $a_{n_k}\in I_k$. Then $|a_{n_k}-x|\le2^{-k}(B-A)\to0$.

\satz{Subsequences keep the limit}{Teilfolgen}
If $a_n\to a$, every subsequence tends to $a$. Hence two different accumulation points mean divergence.
\idea $n_k\ge k$, so $k\ge N$ gives $|a_{n_k}-a|<\eps$.

\satz{Bounded with one accumulation point}{beschränkt, ein Häufungspunkt}
A bounded sequence with exactly one accumulation point $h$ converges to $h$.
\idea Otherwise infinitely many terms stay $\eps_0$ away from $h$. Bolzano--Weierstraß on those gives a second accumulation
point. Unbounded, it fails: $n^{(-1)^n}$.

\satz{Cauchy criterion}{Cauchy-Kriterium}
A sequence in $\R$ converges iff it is a Cauchy sequence.
\idea $\Rightarrow$: $|a_n-a_m|\le|a_n-a|+|a-a_m|<\frac\eps2+\frac\eps2$. $\Leftarrow$: Cauchy $\Rightarrow$ bounded (play
$\eps=1$) $\Rightarrow$ Bolzano--Weierstraß gives $a_{n_k}\to a$ $\Rightarrow$ $|a_n-a|\le|a_n-a_{n_k}|+|a_{n_k}-a|<\eps$ for a
large $k$.

\subsection{Chapter 3}

\satz{Terms of a convergent series tend to 0}{notwendige Bedingung}
\idea $a_n=S_n-S_{n-1}\to S-S=0$. The converse is false: $\sum\frac1n$.

\satz{Cauchy criterion for series}{Cauchy-Kriterium für Reihen}
$\sum a_n$ converges iff $\forall\eps\ \exists N\ \forall n>m\ge N: \bigl|\sum_{k=m+1}^na_k\bigr|<\eps$.
\idea That sum is $S_n-S_m$: this is the Cauchy property of the partial sums.

\satz{Nonnegative terms}{Reihen mit nichtnegativen Gliedern}
If $a_n\ge0$, then $\sum a_n$ converges iff the partial sums are bounded.
\idea The partial sums increase. Monotone convergence.

\satz{The harmonic series diverges}{harmonische Reihe}
\idea Group: $S_{2^k}=1+\frac12+\bigl(\frac13+\frac14\bigr)+\bigl(\frac15+\cdots+\frac18\bigr)+\cdots\ge1+\frac k2$. Each group has
$2^{j-1}$ terms of size at least $2^{-j}$.

\satz{$\sum\frac1{n^\alpha}$ converges for $\alpha>1$}{$p$-Reihe}
\idea The block $2^j\le n<2^{j+1}$ has $2^j$ terms, each at most $2^{-j\alpha}$, so the block is at most $(2^{1-\alpha})^j$. The
partial sums are bounded by the geometric series $\frac1{1-2^{1-\alpha}}$.

\satz{Leibniz criterion}{Leibniz-Kriterium}
$(a_n)$ decreasing, $a_n\ge0$, $a_n\to0$ $\Rightarrow$ $\sum(-1)^na_n$ converges.
\idea Even partial sums decrease, odd ones increase, and $S_{2k+1}\le S_{2k}$. Both converge by monotone convergence, and they
differ by $a_{2k+1}\to0$.

\satz{Absolute convergence implies convergence}{absolute Konvergenz}
\idea $\bigl|\sum_{k=m+1}^na_k\bigr|\le\sum_{k=m+1}^n|a_k|<\eps$. Cauchy criterion.

\satz{Majorant criterion}{Majorantenkriterium}
$|a_n|\le c_n$ for large $n$ and $\sum c_n$ convergent $\Rightarrow$ $\sum a_n$ converges absolutely.
\idea $\sum_{k=m+1}^n|a_k|\le\sum_{k=m+1}^nc_k<\eps$. Cauchy criterion.

\satz{Ratio test}{Quotientenkriterium}
$a_n\ne0$ and $|a_{n+1}/a_n|\le\theta<1$ for $n\ge n_0$ $\Rightarrow$ absolutely convergent.
\idea Induction: $|a_n|\le|a_{n_0}|\theta^{n-n_0}$. A geometric majorant.

\satz{Root test}{Wurzelkriterium}
$\sqrt[n]{|a_n|}\le q<1$ for $n\ge N$ $\Rightarrow$ absolutely convergent; $\sqrt[n]{|a_n|}\ge1$ for infinitely many $n$
$\Rightarrow$ divergent.
\idea $|a_n|\le q^n$: geometric majorant. Otherwise $|a_n|\ge1$ infinitely often: the terms do not tend to $0$.

\satz{Cauchy product}{Cauchy-Produkt}
$\sum a_n$, $\sum b_n$ absolutely convergent $\Rightarrow$ $\sum c_n$, $c_n=\sum_{k=0}^na_kb_{n-k}$, converges absolutely to
$\bigl(\sum a_n\bigr)\bigl(\sum b_n\bigr)$.
\idea $\sum_{n\le N}|c_n|\le\bigl(\sum|a_k|\bigr)\bigl(\sum|b_\ell|\bigr)$, so absolute convergence. $A_NB_N-C_N$ is the sum of
$a_kb_\ell$ over the square $k,\ell\le N$ minus the triangle $k+\ell\le N$; there $k>\frac N2$ or $\ell>\frac N2$, so it is bounded
by tails of $\sum|a_k|$ and $\sum|b_\ell|$, which tend to $0$.

\satz{The exponential series}{Exponentialfunktion}
$\exp x$ converges absolutely for every $x$, and $\exp(x+y)=\exp x\exp y$.
\idea Ratio $\frac{|x|}{n+1}\le\frac12$ once $n+1\ge2|x|$. The functional equation: the Cauchy product's $n$-th term is
$\sum_k\frac{x^ky^{n-k}}{k!(n-k)!}=\frac{(x+y)^n}{n!}$ by the binomial theorem. Consequences: $\exp x\exp(-x)=1$, so $\exp x\ne0$
and $\exp(-x)=\frac1{\exp x}$; $\exp x>0$; $\exp n=e^n$.

\subsection{Chapter 4}

\satz{Sequences and epsilon-delta agree}{Folgen- und $\eps$-$\delta$-Kriterium}
$\lim_{x\to y}f(x)=c$ iff $\forall\eps\ \exists\delta\ \forall x\in D: 0<|x-y|<\delta\Rightarrow|f(x)-c|<\eps$.
\idea From $\eps$-$\delta$ to sequences: a sequence $x_n\to y$ with $x_n\ne y$ eventually lies in the window. Back, by
contradiction: if $\eps_0$ defeats every $\delta$, pick $x_n$ with $0<|x_n-y|<\frac1n$ and $|f(x_n)-c|\ge\eps_0$. Then $x_n\to y$,
but $f(x_n)\not\to c$.

\satz{Continuity is preserved}{Rechenregeln, Verkettung}
Sums, multiples, products, quotients (denominator $\ne0$) and compositions of continuous functions are continuous.
\idea Sequence criterion and the limit laws. Composition: $x_n\to y\Rightarrow f(x_n)\to f(y)\Rightarrow g(f(x_n))\to g(f(y))$.

\satz{$\exp$, $\sin$, $\cos$ are continuous}{Stetigkeit elementarer Funktionen}
\idea $|\exp x-1|\le2|x|$ for $|x|\le1$, and $\exp x=\exp y\exp(x-y)$. For $\sin$, $\cos$: $\frac{\sin h}h\to1$, $\cos h\to1$,
then the addition theorems.

\satz{Intermediate value theorem}{Zwischenwertsatz}
$f\colon[a,b]\to\R$ continuous and $c$ between $f(a)$ and $f(b)$ $\Rightarrow$ some $p\in[a,b]$ has $f(p)=c$.
\idea Reduce to $c=0$, $f(a)<0<f(b)$. Bisect, keeping $f(a_n)\le0\le f(b_n)$. $a_n$ increases, $b_n$ decreases, and
$b_n-a_n=2^{-n}(b-a)$, so both tend to one $p$. Continuity: $f(p)\le0\le f(p)$.

\satz{Extreme value theorem}{Satz vom Maximum und Minimum}
$f\colon[a,b]\to\R$ continuous $\Rightarrow$ $f$ is bounded and attains a maximum and a minimum.
\idea Unbounded would give $|f(x_n)|\ge n$; Bolzano--Weierstraß gives $x_{n_k}\to p\in[a,b]$ (closed!), and $f(x_{n_k})\to f(p)$
is bounded. Contradiction. Maximum: $M=\sup f([a,b])$; pick $f(x_n)>M-\frac1n$; a subsequence tends to some $p$, and $f(p)=M$.

\satz{Heine--Cantor}{Satz von Heine--Cantor}
$f\colon[a,b]\to\R$ continuous $\Rightarrow$ $f$ uniformly continuous.
\idea Otherwise there are $\eps_0$ and $x_n,y_n$ with $|x_n-y_n|<\frac1n$, $|f(x_n)-f(y_n)|\ge\eps_0$. Bolzano--Weierstraß:
$x_{n_k}\to p\in[a,b]$, and then $y_{n_k}\to p$. Continuity at $p$: $f(x_{n_k})-f(y_{n_k})\to0$. Contradiction.

\satz{Continuous extension to the ends}{stetige Fortsetzung an Randpunkte}
$f\colon(a,b)\to\R$ uniformly continuous $\Rightarrow$ a unique continuous $\tilde f\colon[a,b]\to\R$ agrees with $f$ on $(a,b)$.
\idea $x_n\to a$ is Cauchy, so $(f(x_n))$ is Cauchy (uniform continuity) and converges to some $A$. Interleave two such
sequences into one: the limit does not depend on the choice. Set $\tilde f(a)=A$; the same at $b$.

\satz{Approximation by step functions}{Approximation durch Treppenfunktionen}
$f\colon[a,b]\to\R$ continuous, $\eps>0$ $\Rightarrow$ step functions $\varphi\le f\le\psi$ with $\psi-\varphi<\eps$.
\idea Heine--Cantor gives $\delta$. Partition $[a,b]$ into pieces shorter than $\delta$. On each piece $f$ attains a minimum
$m_k$ and a maximum $M_k$ (extreme value theorem), and $M_k-m_k<\eps$. Let $\varphi=m_k$ and $\psi=M_k$ on the pieces.

\satz{Inverse function theorem}{Satz über die Umkehrfunktion}
$I$ an interval, $f\colon I\to\R$ continuous and strictly monotone $\Rightarrow$ $f(I)$ is an interval, $f\colon I\to f(I)$ is
bijective, and $f^{-1}$ is continuous and strictly monotone in the same direction.
\idea Interval: IVT. Injective: strict monotonicity. Continuity of $f^{-1}$ at $y_0=f(x_0)$ (increasing case): with
$y_\pm=f(x_0\pm\eps)$ and $\delta=\min\{y_0-y_-,\ y_+-y_0\}$, any $|y-y_0|<\delta$ lies between $y_-$ and $y_+$, so
$f^{-1}(y)$ lies between $x_0-\eps$ and $x_0+\eps$.

\satz{$\exp$ and $\log$}{Exponentialfunktion und Logarithmus}
$\exp\colon\R\to(0,\infty)$ is continuous, strictly increasing and bijective; $\log=\exp^{-1}$ is continuous, strictly increasing,
and $\log(xy)=\log x+\log y$.
\idea $\exp h>1$ for $h>0$, so $\exp y=\exp x\exp(y-x)>\exp x$ for $y>x$. $\exp n\ge1+n\to\infty$ and $\exp(-n)\to0$; the IVT
gives every $r>0$. For $\log$: apply $\exp$ to both sides and use injectivity.

\satz{The exponential functional equation}{Charakterisierung von $a^x$}
$F\colon\R\to\R$ continuous with $F(x+y)=F(x)F(y)$ $\Rightarrow$ $F\equiv0$, or $F(x)=a^x$ with $a=F(1)>0$.
\idea $F(0)=F(0)^2$. If $F(0)=0$, then $F\equiv0$. Otherwise $F(x)=F(\frac x2)^2\ge0$ and $F(x)F(-x)=1$, so $F>0$. Then $F(n)=a^n$,
$F(\frac pq)^q=a^p$ gives $F(\frac pq)=a^{p/q}$, and density of $\Q$ with continuity gives every real $x$.

% =====================================================================================
\section{Counterexamples}

\opening{``B}{eweisen oder widerlegen''} is the favourite question of every hard examiner, and it punishes the student who starts
proving before checking. Spend sixty seconds trying the usual monsters first. A counterexample is the cheapest proof in
mathematics: one line, and the claim is dead.
\marginkael{questioning}

\begin{panel}{Bestiary}{try these first}{}
\begin{itemize}
  \item $(-1)^n$: bounded, two accumulation points, divergent.
  \item $n^{(-1)^n}$: one accumulation point, unbounded, divergent.
  \item $\sqrt n$ and the harmonic sums $H_n$: steps shrink to $0$, still divergent.
  \item $\sum\frac1n$ and $\sum\frac1{n^2}$: the same ratio limit $1$, opposite fates. $\sum\frac{(-1)^n}n$: convergent, not
        absolutely.
  \item $D$, $1$ on $\Q$ and $0$ elsewhere: continuous nowhere. $xD(x)$: continuous exactly at $0$.
  \item $x^2$ on $\R$ and $\frac1x$ on $(0,1]$: continuous, not uniformly. $\sqrt x$: uniformly continuous, not Lipschitz.
  \item The jump: $0$ for $x<0$, $1$ for $x\ge0$. $(0,1)$ and $\{q\in\Q:q^2<2\}$: sets missing their bound.
\end{itemize}
\end{panel}

\needspace{12\baselineskip}\begin{panel}{False}{numbers and sets}{}
\begin{twocol}{52mm}{Tempting claim}{Killed by}
\row{A bounded set has a maximum.}{$(0,1)$: $\sup=1\notin(0,1)$.}
\row{A bounded set of rationals has a supremum in $\Q$.}{$\{q\in\Q: q^2<2\}$.}
\row{Nested intervals always share a point.}{$(0,\frac1n)$: open. $[n,\infty)$: unbounded.}
\row{$x^2\le y^2\Rightarrow x\le y$.}{$x=1$, $y=-2$. Squares forget signs; compare $|x|$, $|y|$.}
\row{$\mathbb Z$ is a field.}{$2$ has no inverse in $\mathbb Z$.}
\row{$\mathbb C$ can be ordered.}{$i^2=-1$ would be a positive square.}
\end{twocol}
\end{panel}

\needspace{12\baselineskip}\begin{panel}{False}{sequences}{}
\begin{twocol}{52mm}{Tempting claim}{Killed by}
\row{Bounded $\Rightarrow$ convergent.}{$(-1)^n$.}
\row{Monotone $\Rightarrow$ convergent.}{$n$.}
\row{Exactly one accumulation point $\Rightarrow$ convergent.}{$n^{(-1)^n}$. (True if bounded.)}
\row{$|a_{n+1}-a_n|\to0\Rightarrow$ convergent.}{$\sqrt n$; $H_n$.}
\row{$a_nb_n\to0\Rightarrow a_n\to0$ or $b_n\to0$.}{$(1,0,1,0,\dots)$ and $(0,1,0,1,\dots)$.}
\row{$a_n+b_n$ convergent $\Rightarrow$ both convergent.}{$(-1)^n$ and $-(-1)^n$.}
\row{$a_n>0$ for all $n$, $a_n\to a\Rightarrow a>0$.}{$\frac1n\to0$. Limits keep only $\ge$.}
\row{A sequence has finitely many accumulation points.}{An enumeration of $\Q$: every real number is one.}
\row{Solve $L=g(L)$, done.}{$a_{n+1}=2a_n$, $a_1=1$.}
\end{twocol}
\end{panel}

\needspace{12\baselineskip}\begin{panel}{False}{series}{}
\begin{twocol}{52mm}{Tempting claim}{Killed by}
\row{$a_n\to0\Rightarrow\sum a_n$ converges.}{$\sum\frac1n$.}
\row{Convergent $\Rightarrow$ absolutely convergent.}{$\sum\frac{(-1)^n}n$.}
\row{$|a_{n+1}/a_n|<1$ for all $n$ $\Rightarrow$ convergent.}{$\frac1n$: ratio $\frac n{n+1}<1$. You need a fixed $\theta<1$.}
\row{Ratio limit $1$ decides.}{$\sum\frac1n$ diverges, $\sum\frac1{n^2}$ converges.}
\row{Leibniz works without monotonicity.}{$b_{2k}=\frac1k$, $b_{2k+1}=\frac1{k^2}$: the pairs sum to $\frac1k-\frac1{k^2}$,
which diverges.}
\row{Two convergent series have a convergent Cauchy product.}{$\sum\frac{(-1)^n}{\sqrt{n+1}}$ squared (Move~09).}
\row{$\sum a_n$ convergent $\Rightarrow\sum a_n^2$ convergent.}{$a_n=\frac{(-1)^n}{\sqrt n}$.}
\row{$\sum a_n$ convergent, $(b_n)$ bounded $\Rightarrow\sum a_nb_n$ convergent.}{$a_n=\frac{(-1)^n}n$, $b_n=(-1)^n$.}
\end{twocol}
\end{panel}

\needspace{12\baselineskip}\begin{panel}{False}{functions}{}
\begin{twocol}{52mm}{Tempting claim}{Killed by}
\row{Continuous on $(a,b)$ $\Rightarrow$ bounded.}{$\frac1x$ on $(0,1)$.}
\row{Continuous on $[a,b)$ $\Rightarrow$ attains a maximum.}{$x$ on $[0,1)$.}
\row{Continuous $\Rightarrow$ uniformly continuous.}{$x^2$ on $\R$; $\frac1x$ on $(0,1]$.}
\row{Uniformly continuous $\Rightarrow$ bounded.}{$x$ on $\R$.}
\row{Uniformly continuous $\Rightarrow$ Lipschitz.}{$\sqrt x$ on $[0,1]$.}
\row{The IVT holds without continuity.}{The jump at $0$ on $[-1,1]$.}
\row{The IVT holds on $\Q$.}{$x^2-2$ on $[0,2]\cap\Q$.}
\row{$f+g$ continuous $\Rightarrow f,g$ continuous.}{$D$ and $1-D$.}
\row{$|f|$ continuous $\Rightarrow f$ continuous.}{$2D-1$, with values $\pm1$.}
\row{The limit exists at $x_0$ $\Rightarrow$ continuous at $x_0$.}{$0$ for $x\ne0$, $1$ at $0$.}
\row{Injective $\Rightarrow$ monotone.}{$x$ on $[0,1)$, $3-x$ on $[1,2]$. (True for continuous $f$ on an interval.)}
\row{Continuous bijection $\Rightarrow$ continuous inverse.}{$x$ on $[0,1)$, $x-1$ on $[2,3]$, onto $[0,2]$: the inverse jumps at
$1$. The domain is not an interval.}
\end{twocol}
\end{panel}

% =====================================================================================
\section{The exam itself}

\opening{Y}{ou will be} tired at minute seventy, and the professor grades what is on the paper, not what was in your head. These
are the habits that turn understanding into points.

\begin{panel}[unbreakable]{Checklist}{full marks}{volle Punktzahl}
\begin{enumerate}
  \item Quantifiers in the order of the definition; $N$ or $\delta$ named explicitly.
  \item Every inequality valid where you use it. Say where: ``für $n\ge3$'', ``da $|x-1|<1$''.
  \item Every theorem you cite gets its name \emph{and} its hypotheses, checked in writing: IVT (continuous on a closed
        interval), Leibniz (three conditions), Cauchy product (absolute convergence), monotone convergence (monotone and
        bounded), Heine--Cantor ($[a,b]$).
  \item Arrows: $=$ joins objects, $\Rightarrow$ joins statements, $\iff$ only when both directions hold. (Your UE01.)
  \item An equation solved by $\Rightarrow$ gives candidates. Plug them in, check the domain. (Your UE01/7.)
  \item The step that carries the problem is proved, never asserted. (Your UE01/2b.)
  \item Skript tools only: no derivatives, no monotonicity via $f'>0$, no L'Hospital, no Taylor, no integrals, no
        $(1+\frac1n)^n\to e$, no periodicity of $\sin$ or the number $\pi$. Results from the Übungen: safest to reprove them in
        two lines, as this folio does.
\end{enumerate}
\end{panel}

\begin{panel}[unbreakable]{Speed}{ninety minutes}{}
\begin{enumerate}
  \item Minutes 0--3: read every question. Tag each: \textbf{D} definition or statement, \textbf{S} a Skript proof, \textbf{R} a
        move.
  \item D first: the most points per minute. Add one example and one non-example; it costs twenty seconds.
  \item S next, when the Idea line comes at once. Write the idea as a sentence first, then the details.
  \item R: name the move from its ``When'' line. Scratch backwards in the margin, fair copy forwards.
  \item ``Beweisen oder widerlegen'': sixty seconds against the bestiary before you invest in a proof.
  \item Stuck? Write the definition unfolded, all quantifiers, and the first line of the skeleton. Partial credit lives there,
        and so does the next step.
  \item Never polish $N$ or $\delta$. Crude and correct beats tight and late.
  \item Last five minutes: every cited theorem has its hypotheses next to it.
\end{enumerate}
\end{panel}

\begin{panel}[unbreakable]{Phrasebook}{the test is in German}{Formulierungen}
\begin{itemize}
  \item Sei $\eps>0$ beliebig. \ Nach der archimedischen Eigenschaft existiert $N\in\N$ mit $N>\frac3\eps$.
  \item Dann gilt für alle $n\ge N$: \dots \ Somit ist $\lim_{n\to\infty}a_n=\frac12$ gezeigt.
  \item Setze $\delta:=\min\{1,\frac\eps7\}$. Für alle $x\in D$ mit $|x-x_0|<\delta$ gilt \dots
  \item Angenommen, \dots\ Dann \dots, ein Widerspruch. Also gilt \dots
  \item Ohne Einschränkung (o.\,B.\,d.\,A.) sei $x<y$.
  \item Induktionsanfang ($n=1$): \dots \ Induktionsvoraussetzung: Für ein $n\in\N$ gelte \dots \ Induktionsschritt
        ($n\to n+1$): \dots
  \item Die Funktion $h$ ist als Summe stetiger Funktionen stetig auf $[0,1]$. Nach dem Zwischenwertsatz existiert \dots
  \item Die Reihe konvergiert nach dem Majorantenkriterium mit der Majorante $\sum\frac1{n^2}$.
  \item Gegenbeispiel: \dots\ Damit ist die Aussage falsch.
\end{itemize}
\end{panel}

% =====================================================================================
\section{Fourteen days}

\opening{F}{ourteen days} is enough, if every day is a test and none is a reread. About two hours a day. Day~1 is the day you
start; Day~14 is the evening before the test. If the date moves, compress Days~8 and~9.
\aside{Fourteen days. I once watched a star collapse in less, and it had far less to learn.}

\medskip
\begin{plan}
\pday{01}{Parts I and II. Then, on a blank page, write the eight definitions of the negation table with their negations.}
  {all eight, quantifiers in order}
\pday{02}{Part III: the keystone, the terms, Moves 01--03, tasks \ref{t:fac}--\ref{t:comp}. Vorbereitung §1--3 cold.}
  {every task right without the key}
\pday{03}{Part IV: Moves 04--05, tasks \ref{t:epsA}--\ref{t:hp}. Vorbereitung §4--5 cold.}
  {an $\eps$-$N$ proof in under 8 minutes}
\pday{04}{Move 06, Part V: Moves 07--08, tasks \ref{t:rec}--\ref{t:param}. Vorbereitung §6--8 cold.}
  {you name the test before computing}
\pday{05}{Move 09, Part VI: Moves 10--12, tasks \ref{t:leib1}--\ref{t:inf}. Vorbereitung §9--12 cold.}
  {an $\eps$-$\delta$ proof in under 8 minutes}
\pday{06}{Moves 13--16, tasks \ref{t:sq}--\ref{t:inv}. Vorbereitung §13--16 cold.}
  {all sixteen moves named from their ``When'' lines}
\pday{07}{Part VII, chapters 1--2: each theorem's statement and Idea from a blank page. Mark every miss.}
  {misses listed}
\pday{08}{Part VII, chapters 3--4, then Part VIII with the right column covered. Redo Day~7's misses.}
  {each claim killed in one line}
\pday{09}{Your homework, cold: Übung 2/2, 3/4, 3/7, 4/3--4/5, 5/3, 5/6, 5/7. Whatever fails goes to \cmd{/exercise}.}
  {each one written as for the test}
\pday{10}{Mock exam with \cmd{/exam}: ninety minutes, timed, unseen problems, graded strictly.}
  {graded}
\pday{11}{Every mock error into the grimoire, re-attempted the same day; redo the move it belongs to.}
  {every error re-solved}
\pday{12}{Teach-back with \cmd{/teach}: the Cauchy criterion and the IVT, out loud. Then a 20-minute sprint through every term
  of Parts III--VI.}
  {no gaps found}
\pday{13}{A second, shorter mock (60 minutes), or a mixed set of every task you missed.}
  {80\,\% or more}
\pday{14}{Light. Blank-page recall of the triggers and the negation table. Nothing new. Sleep.}
  {asleep before midnight}
\end{plan}

\subsection{The key}
Answers and first moves only. If the first move is not enough, bring the attempt to \cmd{/exercise}: hints come one rung at a
time, and you keep the proof.

\keyline{t:fac}Step: $(n+1)!-1+(n+1)(n+1)!=(n+1)!\,(n+2)-1=(n+2)!-1$.
\keyline{t:fac2}Base $24>16$. Step: $(n+1)!=(n+1)\,n!>(n+1)2^n\ge2\cdot2^n$.
\keyline{t:supmn}$\sup=\max=2$ at $m=n=1$. $\inf=0$, not attained: every element is $>0$; for $\eps$ take $m=n>\frac2\eps$.
\keyline{t:supc}Upper bound $ca\le c\sup A$. Least: pick $a>\sup A-\frac\eps c$.
\keyline{t:bij}$g(x)=(x+1)^2$, strictly increasing; $g^{-1}(y)=\sqrt y-1$, which is $\ge0$ exactly when $y\ge1$.
\keyline{t:comp}Apply $g$ to $f(x_1)=f(x_2)$. For surjectivity, $z=g(f(x))$ is hit by $f(x)$. Example: $M=P=\{1\}$,
$N=\{1,2\}$, $f(1)=1$, $g\equiv1$.
\keyline{t:epsA}$|a_n-3|=\frac{14}{n+4}<\frac{14}n$; $N>\frac{14}\eps$.
\keyline{t:epsB}$|a_n|\le\frac{n+1}{n^2}\le\frac2n$; $N>\frac2\eps$.
\keyline{t:abs}Reverse triangle inequality. The converse fails: $(-1)^n$. It holds when $a=0$.
\keyline{t:sqrt}$a=0$: $\sqrt{a_n}<\eps\iff a_n<\eps^2$. $a>0$: conjugate, $|\sqrt{a_n}-\sqrt a|\le\frac{|a_n-a|}{\sqrt a}$.
\keyline{t:hp}Even $n$: $\to2$. Odd $n$: $\to-1$. Accumulation points $\{-1,2\}$; divergent.
\keyline{t:rec}$L=3$. Bound $0\le a_n<3$. Increasing: $a_{n+1}^2-a_n^2=(3-a_n)(2+a_n)>0$.
\keyline{t:heron}$x_{n+1}-\sqrt a=\frac{(x_n-\sqrt a)^2}{2x_n}\ge0$, then $x_{n+1}-x_n=\frac{a-x_n^2}{2x_n}\le0$ from $n=2$
on. Limit $\sqrt a$.
\keyline{t:cauchy}For $m>n$: $|s_m-s_n|=\sum_{k=n+1}^m2^{-k}<2^{-n}$. Choose $N$ with $2^N>\frac1\eps$ ($2^N\ge1+N$).
\keyline{t:tel}$\frac12$. $\frac1{(2n-1)(2n+1)}=\frac12\bigl(\frac1{2n-1}-\frac1{2n+1}\bigr)$.
\keyline{t:geo}$\frac72=\frac32+2$.
\keyline{t:log}$S_N=\sum_{n=1}^N\bigl(\log(n+1)-\log n\bigr)=\log(N+1)$, unbounded since $\log e^k=k$.
\keyline{t:ratio}Convergent: the ratio is $\frac{(1+1/n)^2}3\le\frac34$ for $n\ge2$.
\keyline{t:root}Convergent: $\frac n{3n-1}\le\frac12$ for all $n\ge1$.
\keyline{t:comp2}Divergent: $\frac1{\sqrt{n(n+1)}}\ge\frac1{n+1}$.
\keyline{t:hard}Divergent: the ratio is $\frac3{(1+1/n)^n}\ge1$, so $a_n\ge a_1=3$: not a null sequence.
\keyline{t:param}Exactly for $|x|<1$. At $|x|\ge1$ the terms do not tend to $0$.
\keyline{t:leib1}Conditionally convergent: Leibniz; and $\frac1{2n+1}\ge\frac1{3n}$ for $n\ge1$.
\keyline{t:leib2}Absolutely convergent: the ratio is $\frac{(1+1/n)^2}2\le\frac89$ for $n\ge3$.
\keyline{t:leib3}Conditionally convergent: $b_n=\frac1{\sqrt{n+1}+\sqrt n}$ decreases to $0$, and $b_n\ge\frac1{2(n+1)}$.
\keyline{t:exp1}$t=\exp x$: $t^2-4t+1=0$, $t=2\pm\sqrt3>0$. So $x=\pm\log(2+\sqrt3)$, as $(2-\sqrt3)(2+\sqrt3)=1$.
\keyline{t:exp2}$1$: divide by $\exp x$; $\exp(-2x)\to0$.
\keyline{t:exp3}$0$: for $x\ge0$, $\exp x\ge\frac{x^3}6$, so $\frac{x^2}{\exp x}\le\frac6x$.
\keyline{t:ed1}$|x^3-1|=|x-1|\,|x^2+x+1|<7|x-1|$ when $|x-1|<1$. $\delta=\min\{1,\frac\eps7\}$.
\keyline{t:ed2}For $x\ne2$ the quotient is $x+2$, and $|x+2-4|=|x-2|$. $\delta=\eps$.
\keyline{t:ed3}$|x-2|<1\Rightarrow x>1\Rightarrow\bigl|\frac1x-\frac12\bigr|=\frac{|x-2|}{2x}<\frac{|x-2|}2$. $\delta=\min\{1,2\eps\}$.
\keyline{t:ext}$f(0)=\frac12$. Conjugate: $f(x)=\frac1{\sqrt{1+x}+1}$.
\keyline{t:inf}For $x\ge3$: $\bigl|f(x)-2\bigr|=\frac{x+6}{x^2-3}\le\frac{3x}{x^2/2}=\frac6x$. $M=\max\{3,\frac6\eps\}$.
\keyline{t:sq}Yes, with $\tilde f(0)=0$: $|\sqrt x\sin\frac1x|\le\sqrt x$.
\keyline{t:dir}At $0$: $|xD(x)|\le|x|$. At $x_0\ne0$: rationals give $q_n\to x_0$, irrationals give $0$; they differ.
\keyline{t:ivt3}$f(-2)<0<f(0)$, $f(1)<0<f(2)$: three disjoint sign changes.
\keyline{t:chord}$h(x)=f(x+\frac12)-f(x)$ on $[0,\frac12]$: $h(0)+h(\frac12)=f(1)-f(0)=0$.
\keyline{t:odd}Factor out the leading term: $p(x)=a_nx^n(1+r(x))$ with $r(x)\to0$ as $|x|\to\infty$. Opposite signs at $\pm R$.
\keyline{t:uc1}$\bigl|\frac1{1+x^2}-\frac1{1+y^2}\bigr|=\frac{|x-y|\,|x+y|}{(1+x^2)(1+y^2)}$, and $\frac{|x|}{1+x^2}\le\frac12$.
\keyline{t:uc2}$(0,1]$: $x_n=\frac1{n+1}$, $y_n=\frac1n$. $[1,\infty)$: $\bigl|\frac1x-\frac1y\bigr|=\frac{|x-y|}{xy}\le|x-y|$.
\keyline{t:uc3}$x_n=n$, $y_n=n+\frac1n$: $\exp y_n-\exp x_n=\exp n\,(\exp\frac1n-1)\ge\frac{\exp n}n>1$.
\keyline{t:inv}For $0\le x<y$: $f(y)-f(x)=\frac{(y-x)(x+y+xy)}{(1+x)(1+y)}>0$. $f(n)\ge\frac n2$, image $[0,\infty)$.
$f^{-1}(y)=\frac{y+\sqrt{y^2+4y}}2$.

\begin{triggers}
\trigger{$|a_n-a|$ is a fraction}{one fraction; numerator up, denominator down, to $\frac Cn$}
\trigger{A minus sign in a denominator}{keep half: $n^2-5\ge\frac{n^2}2$ for $n\ge4$}
\trigger{A difference of roots}{the conjugate}
\trigger{Supremum or infimum}{upper bound, then an element beyond $s-\eps$; Archimedes picks the index}
\trigger{A recursive sequence}{bound by induction, factor $a_{n+1}-a_n$, then $L=g(L)$ and the bound picks the root}
\trigger{Signs that oscillate}{split the indices by the period; two limits mean divergence}
\trigger{A product in a series denominator}{partial fractions, then telescope}
\trigger{$n!$ or $c^n$ in a series}{the ratio test with one fixed $\theta<1$}
\trigger{An alternating series}{$\sum|a_n|$ first; then Leibniz, monotonicity by cross-multiplying}
\trigger{$\eps$-$\delta$ with a product}{restrict $\delta\le1$, bound the other factor}
\trigger{A glue point or a parameter}{one-sided limits equal to the value}
\trigger{``Genau eine Lösung''}{IVT for existence, strict monotonicity for uniqueness}
\trigger{Not uniformly continuous?}{two sequences, closer and closer, where $f$ is steepest}
\trigger{An inverse function}{continuous, strictly monotone, image by the IVT, then the theorem}
\trigger{``Beweisen oder widerlegen''}{sixty seconds against the bestiary first}
\end{triggers}

\begin{deck}
\card{Def}{The supremum $s$ of $M$, in its $\eps$-form.}{$s$ is an upper bound of $M$, and $\forall\eps>0\ \exists x\in M: x>s-\eps$.}
\card{Def}{Uniform continuity, and its negation.}{$\forall\eps\,\exists\delta\,\forall x,y: |x-y|<\delta\Rightarrow|f(x)-f(y)|<\eps$;
fails: $\exists\eps\,\forall\delta\,\exists x,y: |x-y|<\delta$, $|f(x)-f(y)|\ge\eps$.}
\card{Why}{Why is $L=g(L)$ not a proof that $a_{n+1}=g(a_n)$ converges to $L$?}{It assumes the limit exists: $a_{n+1}=2a_n$
gives $L=0$ and diverges.}
\card{Trigger}{Continuous on an open or unbounded set; is it uniformly continuous? Reach for?}{Two sequences with
$|x_n-y_n|\to0$ and $|f(x_n)-f(y_n)|\ge\eps_0$, where $f$ is steepest.}
\card{Trigger}{$n!$ or $c^n$ in a series. Reach for?}{The ratio test with a fixed $\theta<1$; Bernoulli gives
$(1+\frac1n)^n\ge2$.}
\end{deck}

\end{document}
